\documentclass[11pt]{amsart}
\usepackage{amsmath,amssymb,amsthm}
\usepackage{xcolor}
\usepackage[colorlinks=true,linkcolor=blue,citecolor=blue]{hyperref}
\usepackage{microtype}

\numberwithin{equation}{section}
\theoremstyle{plain}
\newtheorem{theorem}{Theorem}[section]
\newtheorem{proposition}[theorem]{Proposition}
\newtheorem{lemma}[theorem]{Lemma}
\newtheorem{corollary}[theorem]{Corollary}
\newtheorem{conjecture}[theorem]{Conjecture}
\theoremstyle{definition}
\newtheorem{definition}[theorem]{Definition}
\newtheorem{convention}[theorem]{Convention}
\newtheorem{example}[theorem]{Example}
\theoremstyle{remark}
\newtheorem{remark}[theorem]{Remark}
\newtheorem{question}[theorem]{Question}

\newcommand{\cA}{\mathcal{A}}
\newcommand{\cB}{\mathcal{B}}
\newcommand{\cC}{\mathcal{C}}
\newcommand{\cM}{\mathcal{M}}
\newcommand{\cW}{\mathcal{W}}
\newcommand{\cZ}{\mathcal{Z}}
\newcommand{\cL}{\mathcal{L}}
\newcommand{\bC}{\mathbb{C}}
\newcommand{\bQ}{\mathbb{Q}}
\newcommand{\bR}{\mathbb{R}}
\newcommand{\bZ}{\mathbb{Z}}
\newcommand{\fg}{\mathfrak{g}}
\newcommand{\gmod}{\mathfrak{g}^{\mathrm{mod}}}
\newcommand{\Ran}{\mathrm{Ran}}
\newcommand{\FM}{\mathrm{FM}}
\newcommand{\Conf}{\mathrm{Conf}}
\newcommand{\barB}{\overline{B}}
\newcommand{\Vir}{\mathrm{Vir}}
\newcommand{\ChirHoch}{\mathrm{ChirHoch}}
\newcommand{\MC}{\mathrm{MC}}
\newcommand{\Def}{\mathrm{Def}}
\newcommand{\dvee}{h^{\vee}}
\newcommand{\Dfib}{d_{\mathrm{fib}}}
\newcommand{\bE}{\mathbb{E}}
\newcommand{\Mbar}{\overline{\cM}}
\newcommand{\kdual}{{}^{!}}
\newcommand{\Ob}{\mathrm{Obs}}

\title{Modular Koszul duality for chiral algebras}
\author{Raeez Lorgat}
\date{July 2026}
\subjclass[2020]{17B69, 18M70, 14H10, 81T40, 81R10, 11F46}

\begin{document}

\begin{abstract}
The product of two chiral fields on an algebraic curve diverges where the fields
collide. Among differential forms with a pole along the diagonal, exactly one has
a residue independent of every auxiliary choice: the logarithmic form
$\eta = d\log(z_1-z_2)$. We show that $\eta$ generates a differential on the
Fulton--MacPherson compactification of configuration space, that the vanishing of
its square is equivalent to the Borcherds identity, and that the resulting bar
construction closes, over the moduli space of curves, into a single
Maurer--Cartan element $\Theta_{\cA}$ in a modular convolution Lie algebra.
The equation $D\Theta_{\cA}+\tfrac12[\Theta_{\cA},\Theta_{\cA}]=0$ holds
identically: it is the algebraic form of
$\partial^2=0$ on the Deligne--Mumford boundary, and it holds identically, term by term.
Every invariant constructed in this paper is a projection of $\Theta_{\cA}$:
the modular characteristic $\kappa$ at degree two, the classical $r$-matrix as
the binary collision residue, an obstruction tower whose generating function
satisfies an algebraic equation of degree two, and a genus expansion whose scalar
part is $\kappa$ times the $\widehat{A}$-genus. The ordered form of the
construction remembers strictly more than its symmetric quotient; the discrepancy
is a quantum group. At positive genus the construction acquires a curvature,
$\kappa$ times the first Chern class of the Hodge bundle, and this single
scalar is simultaneously the conformal anomaly of two-dimensional
physics and the coefficient of the Hodge classes at every genus.
\end{abstract}

\maketitle
\tableofcontents

\section{Introduction}

\subsection{The seed}

Let $X$ be a smooth algebraic curve over $\bC$ and let $a(z)$, $b(w)$ be chiral
fields on $X$. Their product diverges as $z\to w$. Any regularisation of this
divergence requires a $1$-form on the configuration space of two points with a
pole along the diagonal, and among such forms exactly one has a residue that does
not depend on a coordinate, a metric, or a level:
\begin{equation}\label{eq:eta}
\eta_{12} \;=\; d\log(z_1-z_2)\;=\;\frac{dz_1-dz_2}{z_1-z_2}.
\end{equation}
Three facts about $\eta$ generate everything in this paper.

First, consistency at three points forces a relation. On $\Conf_3(\bC)$,
\begin{equation}\label{eq:arnold}
\eta_{12}\wedge\eta_{23}+\eta_{23}\wedge\eta_{31}+\eta_{31}\wedge\eta_{12}=0,
\end{equation}
which is the partial-fractions identity in disguise; Arnold proved that this one
relation generates the full ideal of relations in
$H^{*}(\Conf_n(\bC))$ \cite{Arnold}. Second, extraction of residues against
$\eta$ converts the operator product expansion into a differential, and the
vanishing of the square of that differential is exactly the associativity axiom
of a chiral algebra. Third, $\eta$ fails to be invariant under a change of
coordinate $z\mapsto w(z)$ by the cocycle $d\log w'(z)$; this failure is the
central extension that the Virasoro algebra measures. One $1$-form, one
relation, one extraction: everything that follows is forced.

\subsection{The tower}

The objects of the theory arrange themselves in a tower, and the paper
climbs it:
\[
\begin{array}{ccccc}
\text{factorization category} &\longrightarrow& \cA=\operatorname{End}(b)
&\longrightarrow& \barB_X(\cA)\\[2pt]
&&\downarrow&&\\[2pt]
\text{modular scalars} &\longleftarrow& \text{line operators}
&\longleftarrow& \cZ^{\mathrm{der}}_{\mathrm{ch}}(\cA)
\end{array}
\]
At the base sits a factorization category on the curve with a chosen
vacuum; its endomorphism algebra is a chiral algebra; the bar
construction carries the algebra to a coalgebra on compactified
configurations; the centre is a three-dimensional object; its module
categories carry the line operators; at the top sit the scalars ---
characters, Verlinde numbers, modular forms. Each floor is recovered
from the next by a reconstruction theorem: Morita theory at the base,
bar--cobar inversion above it, terminality of the centre above that,
the modular trace at the top. One element, the Maurer--Cartan element
$\Theta_{\cA}$ of Section~\ref{sec:genus}, generates every floor; one
involution, Koszul duality, acts on all of them.

\subsection{The results}

Sections~\ref{sec:bar}--\ref{sec:koszul} construct the bar complex of a chiral
algebra geometrically, as the space of logarithmic forms on Fulton--MacPherson
compactifications, and prove that the construction is a selection principle:
the differential squares to zero if and only if the input satisfies the
Borcherds identity (Theorem~\ref{thm:selection}). A companion rigidity
result sharpens this: nilpotence consumes the full differential, level
term included (Proposition~\ref{prop:bracket-fails}); it is a
cancellation between geometry and algebra, finer than any Jacobi
identity.

Section~\ref{sec:koszul} separates three operations: reconstruction
(the bar--cobar counit), Koszul duality (Verdier duality on the Ran
space applied to the bar coalgebra), and the derived centre --- three
functors with three targets, related exactly by the theorems that
relate them. The main recognition result is a
Poincar\'e--Birkhoff--Witt criterion (Theorem~\ref{thm:pbw})
establishing Koszulness for the universal affine, Virasoro, $\cW$,
free-field and lattice algebras, and locating its exact boundary:
the Ising model and the simple admissible quotients
$L_k(\mathfrak{sl}_3)$ at denominator $q\geq 3$ lie outside it, each
for an identified structural reason.

Section~\ref{sec:ordered} proves that the ordered bar complex is the primitive
object and the symmetric one its shadow. The kernel of the symmetrisation map
has dimension $d^{\,n}-\binom{n+d-1}{d-1}$ in arity $n$
(Theorem~\ref{thm:reynolds}), and what the kernel carries is precisely the
quantum-group data: the $R$-matrix at degree two, the Drinfeld associator at
degree three, the full braid monodromy beyond. Koszul duality acts on the
ordered theory by inverting the noncommutativity datum, $R\mapsto R^{-1}$,
realised geometrically as orientation reversal of collision cycles
(Theorem~\ref{thm:inversion-principle}); for the Yang $R$-matrix this gives
$Y_{\hbar}(\mathfrak{sl}_2)^{!}\simeq Y_{-\hbar}(\mathfrak{sl}_2)$
(Theorem~\ref{thm:yangian}).

Section~\ref{sec:genus} passes to positive genus. The unique replacement for
$z_1-z_2$ is the prime form; the corrected propagator acquires a
non-holomorphic term, and the fibrewise bar differential curves:
\begin{equation}\label{eq:curvature-intro}
\Dfib^{2}\;=\;\kappa(\cA)\cdot\omega_g,
\end{equation}
where $\omega_g$ is the first Chern class of the Hodge bundle and
$\kappa(\cA)$ is a scalar computed by a two-channel residue calculation
(Theorem~\ref{thm:curvature}). Flatness is restored by the Gauss--Manin
connection, and since the boundary of the boundary of $\Mbar_{g,n}$ is empty,
the total differential $D_{\cA}$ satisfies $D_{\cA}^{2}=0$ unconditionally. Consequently $\Theta_{\cA}:=D_{\cA}-d_0$ is a Maurer--Cartan element
in the modular convolution Lie algebra (Theorem~\ref{thm:mc}); this is the
central object of the paper. For the algebras of central
interest the completion is intrinsic --- an explicit witness forces it
--- and the correct ambient is the coderived category of the second
kind (Theorem~\ref{thm:coderived}).

Section~\ref{sec:shadow} computes the obstruction tower of $\Theta_{\cA}$
restricted to a primary line. The generating function $H(t)=\sum_r rS_rt^r$
of the tower satisfies
\begin{equation}\label{eq:riccati-intro}
H(t)^{2}\;=\;t^{4}\,Q(t),\qquad
Q(t)=(2\kappa+3S_3t)^{2}+2\Delta\,t^{2},\qquad \Delta=8\kappa S_4,
\end{equation}
an algebraic equation of degree two determined by three numbers
(Theorem~\ref{thm:riccati}). The possible depths of the tower are $2$, $3$,
$4$ or $\infty$, and depth $3$ on a single line is impossible
(Theorem~\ref{thm:depth-gap}); this dichotomy is the origin of the
classification of chiral algebras into four classes with a fifth critical
stratum (Theorem~\ref{thm:classification}).

Section~\ref{sec:conductor} proves the complementarity law: the cohomology of
the centre local system over $\Mbar_g$ splits into contributions of $\cA$ and
of its Koszul dual, and the scalar trace of the splitting is the sum
$\varkappa=\kappa(\cA)+\kappa(\cA^{!})$. For the affine and free-field
families $\varkappa=0$; for the Virasoro family $\varkappa=13$, equivalently
$c+c^{!}=26$ --- the ghost budget of the bosonic string arising here as
arithmetic (Theorem~\ref{thm:complementarity}).

Section~\ref{sec:landscape} works the theory family by family ---
Heisenberg, affine, level one, symmetric orbifolds, $\beta\gamma$ ---
until every clause of the classification has a computed witness, and
Section~\ref{sec:DS} crosses the landscape by Drinfeld--Sokolov
reduction: gauge theory becomes gravity, class $\mathsf{L}$ becomes
class $\mathsf{M}$, and the $\cW$-algebras acquire their quantum
groups, including the chiral bialgebra of $\cW_{1+\infty}$ with its
computed antipode obstruction. Section~\ref{sec:moduli} assembles the
recognition theory into a moduli problem over the torsor of
associators and extracts the fingerprint: five numbers that determine
a landscape algebra up to the braiding.

Section~\ref{sec:conductor} proves the complementarity law, the
quantisation law $\hbar^{2}\varkappa=-1$, and the self-dual points.
Section~\ref{sec:genus-expansion} computes the genus expansion. Its scalar
part is a universal multiple of the Faber--Pandharipande numbers,
\begin{equation}
\sum_{g\geq1}F^{\mathrm{sc}}_{g}\,x^{2g}
=\kappa(\cA)\Bigl(\frac{x/2}{\sin(x/2)}-1\Bigr),
\end{equation}
the Wick-rotated $\widehat{A}$-genus (Theorem~\ref{thm:genus-expansion}).
The scalar formula is exact for algebras with a single generator weight and
fails beyond: for $\cW_3$ at genus two the correction is $(c+204)/16c$
(Theorem~\ref{thm:multiweight}), and the corrections dominate the scalar part
at large genus. The genus tower converges; the corresponding string expansion
diverges factorially. The shadow is the convergent projection of a divergent
theory.

Section~\ref{sec:gaudin} reads the ordered theory spectrally --- the
chain from the Maurer--Cartan element through the $r$-matrix and the
transfer matrix to the Baxter operator, with the Bethe ansatz as its
chain-level refinement. Section~\ref{sec:blocks} descends to
conformal blocks: the Verlinde formula from chiral homology, the
Dedekind eta function as a Fredholm determinant, the Kac--Wakimoto
characters from the bar resolution, and the critical level, where the
centre becomes the algebra of functions on opers.
Section~\ref{sec:genus-two} opens the genus-two window: the level
enters the curvature, the double replaces the Yangian, the Igusa form
enters the sewing kernel.

Section~\ref{sec:bulk} identifies the closed sector. The derived chiral
centre, acting on the boundary algebra through the two-coloured Swiss-cheese
operad, is the terminal local open--closed pair
(Theorem~\ref{thm:terminality}): every local closed sector acting on $\cA$
factors uniquely through it. This is holography as a universal property, and
it is sharp: a physical bulk theory is specified by the algebraic data
together with one further datum, an open--closed comparison map, and
Proposition~\ref{prop:separation} shows the datum is irreducible.

Section~\ref{sec:arithmetic} exhibits the two arithmetic channels of the
tower and their asymmetry: the Virasoro tower is the iterated self-extension
of a single Kummer motive --- infinite depth, no transcendence --- while for
lattice algebras the tower resolves an Epstein zeta function one Hecke
eigenform per degree, so that the quartic obstruction of the Leech lattice
algebra is the Ramanujan $L$-function. Kummer-irregular primes enter the
Verlinde polynomials through Bernoulli numbers (Theorem~\ref{thm:irregular})
and provably avoid the Virasoro tower in the computed range.

Section~\ref{sec:E3} completes the local picture --- the derived
centre of the affine algebra is the observable algebra of quantum
Chern--Simons theory on the disc, and the operadic circle
$E_3\to E_2\to E_1\to E_2\to E_3$ closes --- and
Section~\ref{sec:CY} opens the geometric source: the Atiyah--Connes
class of a Calabi--Yau category, the spine theorem carrying it to the
first bar--cobar obstruction of the associated chiral algebra, and
the quantum groups of $K3$. Section~\ref{sec:moonshine} reaches the
automorphic boundary, where the strictification mechanism of the
ordered theory locates the imaginary roots.
Section~\ref{sec:physics} records the physics the formalism derives,
each statement carrying its comparison map, and
Section~\ref{sec:question} states the one problem to which the paper
reduces.Every result stated in this paper is proved in this paper, at the
scope its statement declares; the monograph \cite{Vol1} develops the
surrounding theory at length.

\subsection{Relation to the literature}

The objects assembled here have distinguished ancestries. Chiral
algebras and the Ran space are Beilinson--Drinfeld \cite{BD}; the
categorical bar--cobar equivalence in that setting is Francis--Gaitsgory
\cite{FG}; factorization algebras as quantum field theory are
Costello--Gwilliam \cite{CG}. The compactifications are
Fulton--MacPherson \cite{FM}, the relations Arnold \cite{Arnold}, the
modular operads and their Feynman transform Getzler--Kapranov
\cite{GK}, the curved coalgebras and the second-kind derived
categories Positselski \cite{Pos}. The Hodge-integral evaluations are
Mumford and Faber--Pandharipande \cite{Mum,FP}; the prime form is
Fay \cite{Fay}; the $r$-matrices and quantum groups are Drinfeld
\cite{Drinfeld}, their chiral quantization Etingof--Kazhdan \cite{EK},
their geometric action Maulik--Okounkov \cite{MO} and
Costello--Witten--Yamazaki \cite{CWY}. What is new is the single
construction that holds them together: the geometric bar complex whose
nilpotence is the Borcherds identity, its curvature over moduli, the
Maurer--Cartan element it defines, and the computation of that
element's projections across the entire standard landscape.

\subsection*{Conventions}

We work over $\bC$. Complexes are cohomologically graded; differentials have
degree $+1$. For a graded vector space $V$, $sV$ and $s^{-1}V$ denote
suspension and desuspension. A chiral algebra on $X$ is a right
$\mathcal{D}_X$-module $\cA$ with a chiral product
$\mu\colon j_{*}j^{*}(\cA\boxtimes\cA)\to\Delta_{!}\cA$ satisfying locality,
unitality and the Borcherds identity, in the sense of Beilinson and
Drinfeld \cite{BD}; augmented means equipped with a splitting
$\varepsilon\colon\cA\to\omega_X$, with augmentation ideal
$\bar\cA=\ker\varepsilon$. Locally, for fields $a,b$ the product is encoded by
the operator product expansion
$a(z)b(w)\sim\sum_{n\geq0}(a_{(n)}b)(w)/(z-w)^{n+1}$.

\section{The bar construction on configuration spaces}\label{sec:bar}

\subsection{Compactified configurations}

Let $\Conf_n(X)\subset X^{n}$ be the space of $n$ distinct labelled points and
let $\FM_n(X)$ be its Fulton--MacPherson compactification \cite{FM}: the
iterated blow-up of $X^{n}$ along the partial diagonals in increasing order of
dimension. It is a smooth manifold with corners whose boundary is a normal
crossing divisor $D=\bigcup_{S}D_S$ indexed by subsets $S$ with $|S|\geq2$;
strata are indexed by nested families of subsets, and the open part of $D_S$
fibres over $\FM_{n-|S|+1}(X)$ with fibre a screen: the space of $|S|$ points
on the tangent line modulo affine transformations. On a curve the pairwise
diagonals are already divisors, so the screen of a pair is a point; the first
genuine blow-up occurs at triple collisions.

The form \eqref{eq:eta} extends to $\FM_n(X)$ with logarithmic singularities
along $D$; its residue along $D_S$ equals $1$ if $\{i,j\}\subseteq S$ and $0$
otherwise. This residue is the only memory the compactification keeps of a
collision, and it is independent of every choice. The Arnold
relation \eqref{eq:arnold} is the statement that the forms $\eta_{ij}$ know
about triple collisions; it is the entire codimension-two constraint on a
curve.

\subsection{The ordered bar complex}

\begin{definition}\label{def:bar}
Let $(\cA,\varepsilon)$ be an augmented chiral algebra on $X$. The
\emph{ordered bar complex} of $\cA$ is the cofree conilpotent tensor
coalgebra
\begin{equation}\label{eq:barB}
\barB^{\mathrm{ord}}_{X}(\cA)\;=\;T^{c}(s^{-1}\bar\cA)
\;=\;\bigoplus_{n\geq1}\,(s^{-1}\bar\cA)^{\otimes n}
\otimes\Gamma\bigl(\FM^{\mathrm{ord}}_{n}(X),\,\Omega^{n-1}_{\log}\bigr),
\end{equation}
with the deconcatenation coproduct
$\Delta[a_1|\cdots|a_n]=\sum_{p=0}^{n}[a_1|\cdots|a_p]\otimes
[a_{p+1}|\cdots|a_n]$ and the differential
\begin{equation}\label{eq:bardiff}
d\;=\;d_{\mathrm{int}}+d_{\mathrm{res}}+d_{\mathrm{dR}},
\end{equation}
where $d_{\mathrm{int}}$ is the internal differential of $\cA$,
$d_{\mathrm{dR}}$ the de Rham differential of logarithmic forms, and
\begin{equation}\label{eq:dres}
d_{\mathrm{res}}\bigl([a_1|\cdots|a_n]\otimes\alpha\bigr)
=\sum_{i=1}^{n-1}\sum_{r\geq0}
(-1)^{\epsilon_i}\,
[a_1|\cdots|a_{i}{}_{(r)}a_{i+1}|\cdots|a_n]
\otimes\mathrm{Res}_{D_{i,i+1}}\alpha,
\end{equation}
with $\epsilon_i=\sum_{q<i}(|a_q|-1)+|a_i|$. The sum in \eqref{eq:dres} runs
over \emph{all} polar modes of the operator product: the logarithmic kernel
absorbs one order of pole, so a pole of order $r+1$ in the OPE contributes the
mode $a_{(r)}b$.
\end{definition}

The differential extracts consecutive collisions only; this is what the word
\emph{ordered} means, and the boundary faces that appear are exactly the
$n-1$ consecutive-pair divisors, the facets of the Stasheff associahedron.
The homotopy structure is integral-geometric: realising the associahedron
$K_n$ as the real ordered compactification inside the complex one, the
transferred $A_\infty$ operations are the integrals
$m_n=\int_{K_n}Y(a_1,z_1)\cdots Y(a_n,z_n)\,\omega_n$, the Stasheff
identities are the cancellation of the codimension-one facets, and the
pentagon first appears at codimension three --- coherence is Stokes'
theorem, applied to a polytope.
The symmetric bar complex $\barB^{\Sigma}_{X}(\cA)$ is the
$\Sigma_n$-coinvariant quotient, a factorization coalgebra on the Ran space
of $X$; the relation between the two is the subject of
Section~\ref{sec:ordered}.

\subsection{The selection principle}

\begin{theorem}[Selection principle]\label{thm:selection}
Let $\cA$ be an augmented $\mathcal{D}_X$-module equipped with binary
operations $a_{(r)}b$ of all orders. The operator $d$
of \eqref{eq:bardiff} satisfies $d^{2}=0$ if and only if the operations
satisfy the Borcherds identity. The two constraints decouple by stratum:
on codimension-two boundary strata with disjoint supports, residues
anticommute for orientation reasons; on triple-collision strata the cyclic
sum of iterated residues vanishes if and only if the Borcherds identity
holds, the geometric coefficient being supplied by the Arnold
relation \eqref{eq:arnold}.
\end{theorem}

\begin{proof}
Stokes' theorem on the manifold with corners $\FM_n(X)$ reduces
$d^{2}$ to a sum over codimension-two boundary strata, and the strata
sort into three types. \emph{Disjoint pairs}: a stratum
$D_{ij}\cap D_{kl}$ with $\{i,j\}\cap\{k,l\}=\varnothing$ is reached
through its two codimension-one parents in either order; the two
iterated residues carry opposite orientation signs, from the
anticommutation $e_{ij}\wedge e_{kl}=-e_{kl}\wedge e_{ij}$ of the
conormal directions, and cancel pairwise. \emph{Repeated pairs}: a
single divisor contributes one Poincar\'e residue only, the first
extraction consuming the unique logarithmic normal factor, so no
second-order term arises. \emph{Triple collisions}: the stratum over
$D_{ijk}$ is approached through its three parents
$D_{ij}, D_{jk}, D_{ik}$; the three iterated residues assemble, by
the Arnold relation \eqref{eq:arnold} applied to the form factors,
into the cyclic sum of compositions
$(a_{(m)}b)_{(n)}c$ over all admissible mode orders, and the
vanishing of this sum for all modes and all $c$ is precisely the
Borcherds identity. The geometry supplies the coefficient, the
algebra the cancellation; each is used exactly once. Conversely, a
triple of modes violating the identity produces a chain with
$d^{2}\neq0$ supported on the corresponding stratum, so the
implication is an equivalence; the orientation bookkeeping is
completed in \cite{Vol1}.
\end{proof}

Theorem~\ref{thm:selection} says the bar construction \emph{recognises} chiral
algebras: the Borcherds identity is the exact condition the geometry
accepts. Configuration-space geometry
accepts exactly the algebras satisfying the Borcherds identity. The
following negative result shows how little of the differential is dispensable.

\begin{proposition}\label{prop:bracket-fails}
Let $\cA=V_k(\mathfrak{sl}_2)$ and let $d_{\mathrm{br}}$ denote the part
of $d_{\mathrm{res}}$ retaining only the simple-pole modes $a_{(0)}b$.
Then $d_{\mathrm{br}}^{2}\neq0$, for every one of the $2^{11}$ admissible
assignments of signs to the terms of $d_{\mathrm{br}}$. Nilpotence of the
full differential requires the double-pole (level) term.
\end{proposition}

Nilpotence consumes the level term: the central extension is
structure, spent in the proof of $d^{2}=0$.

\subsection{First computations}

Write $B^{n}$ for the bar complex in bar degree $n$ and $H^{n}$ for its
cohomology, on $X=\mathbb{P}^{1}$ at a fixed point of the augmentation.

\begin{example}[Heisenberg]\label{ex:heisenberg}
For the rank-one Heisenberg algebra $\mathcal{H}_k$, generated by $\alpha$
with $\alpha_{(1)}\alpha=k$, the residue differential vanishes on words in
$\alpha$ except through the double pole, which produces a scalar: the bar
complex is quadratic, the cohomology is
$(\bC,\,0,\,\bC\!\cdot\!c_k,\,0,\dots)$, and the class in degree two is the
central extension of the loop algebra. The dual object is a curved
symmetric coalgebra with curvature $-k\omega$; it is \emph{not} the
Heisenberg algebra at level $-k$.
\end{example}

\begin{example}[Free fermion]
For the free fermion the cyclic group $\bZ/3$ acts on the two-dimensional
space $\Omega^2_{\log}(\FM_3)$ with eigenvalues the primitive cube roots of
unity; total antisymmetry of $\psi\otimes\psi\otimes\psi$ forces invariance;
the intersection is zero and the bar cohomology collapses to $\bC$.
\end{example}

\begin{example}[Affine Kac--Moody]\label{ex:km}
For $\widehat{\fg}_k$ the bar chain space in degree $n$ has dimension
$(\dim\fg)^{n}\,(n-1)!$ --- the factor $(n-1)!$ is the rank of
$\Omega^{n-1}_{\log}$, the Arnold--Orlik--Solomon count --- and the level
enters only the differential, not the chain space. For
$\widehat{\mathfrak{sl}}_2$ the cohomology has dimension $2n+1$ in degree
$n$, with generating function $x(3-x)/(1-x)^2$. The cohomology is
Chevalley--Eilenberg cohomology of the loop algebra corrected by the
Orlik--Solomon factor, and the factor is exact: dimension $2n+1$ in
degree $n$.
\end{example}

\begin{example}[Virasoro]\label{ex:vir}
Every presentation of the Virasoro algebra carries a quartic pole: the operator
product of $T$ with itself has a fourth-order pole and $T$ appears on both
sides. Its bar cohomology has dimensions given by first differences of the
Motzkin numbers, with generating function
$4x/\bigl(1-x+\sqrt{1-2x-3x^2}\,\bigr)^{2}$. In degree two, the vacuum
term $(c/2)[T]\otimes(\eta_{12}+\eta_{13}+\eta_{23})$ survives:
the Arnold relation is quadratic, not linear, and the surviving class is
the germ of the curvature of Section~\ref{sec:genus}.
\end{example}

\subsection{Global sections and boundaries}

Two global statements complete the local picture. First, the homology of
the geometric bar complex computes factorization homology: for holonomic
$\cA$, Weiss excision identifies
$H_{*}\bigl(\barB^{\mathrm{geom}}(\cA)\bigr)\cong\int_{X}\cA$, so the bar
construction is a chain-level model for the global invariant, with the
coradical filtration realised geometrically by the depth of collision
trees. Second, on a curve with marked points and tangent directions the
correct compactification is the bordered one: real oriented blow-up
replaces punctures by boundary circles, and the codimension-one boundary
of the fixed-curve moduli has exactly three types of face --- interior
collision, carrying the screens above; consecutive boundary collision,
carrying Stasheff associahedra and extracting $A_\infty$ operations on
boundary modules; and mixed bubbling, carrying the Swiss-cheese operad
and extracting bulk-to-boundary maps. The fourth type of face, nodal
degeneration of the curve itself, exists only in the relative situation
over $\Mbar_{g,n}$ and is the subject of Section~\ref{sec:genus}. The
three fixed-curve face types are the complete geometric source of the
open--closed structure of Section~\ref{sec:bulk}, complete as it
stands.

\subsection{Modules}

Modules enter through pointed configurations: the module bar complex
of an $\cA$-module $M$ lives on the compactification with one marked
point, is a comodule over $\barB_X(\cA)$ by deconcatenation at the
marked point, and its degree-zero cohomology with several insertions
is the space of derived coinvariants whose comparison with conformal
blocks occupies Section~\ref{sec:blocks}. Every duality of this paper
has a module form, obtained by marking a point in its proof.

\section{Reconstruction, duality, and the centre}\label{sec:koszul}

\subsection{Three functors}

\begin{definition}[Five objects]\label{conv:firewall}
Five objects are associated to an augmented chiral algebra $\cA$, in
five different categories:
\begin{enumerate}
\item the algebra $\cA$ itself;
\item the bar coalgebra $\barB_X(\cA)$;
\item for a quadratic presentation $\cA=T_X(V)/(R)$, the quadratic
coalgebra $\cA^{i}=C_X(s^{-1}V,\,s^{-2}R)$ with its canonical comparison
$q_{\cA}\colon\cA^{i}\to\barB_X(\cA)$;
\item the Verdier algebra
$K_X(\cA)=\mathbb{D}_{\Ran}\,\barB_X(\cA)$, the Verdier dual of the bar
coalgebra as a holonomic object on the Ran space, together with a chosen
pair map $\nu_{\cA}\colon K_X(\cA)\to\cA^{!}=(\cA^{i})^{\vee}$;
\item the derived chiral centre
$\cZ^{\mathrm{der}}_{\mathrm{ch}}(\cA)
=R\mathrm{Hom}_{\cA\otimes\cA^{\mathrm{op}}}(\cA,\cA)$.
\end{enumerate}
The counit $\Omega_X\barB_X(\cA)\to\cA$ of the bar--cobar adjunction is
\emph{reconstruction}: a covariant operation with value $\cA$. Koszul
duality is the contravariant operation (4). The centre is the third
operation (5). An equivalence between the covariant and the
contravariant outputs is precisely a self-duality datum on $\cA$.
\end{definition}

The adjunction itself is classical in the enhanced setting: for the reduced
associative operad in the pro-nilpotent Ran ambient of Francis and Gaitsgory
the bar and cobar functors are mutually inverse equivalences between
augmented algebras and conilpotent coalgebras \cite{FG}. What requires proof
in the present geometric model is, first, that the chain-level realisation
on Fulton--MacPherson spaces computes this equivalence --- this holds on the
completed finite-type locus, by a Mittag--Leffler argument recalled in
Section~\ref{sec:genus} --- and, second, when the quadratic comparison
$q_{\cA}$ is a quasi-isomorphism. The latter is Koszulness.

\subsection{The recognition theorem}

\begin{theorem}[PBW criterion]\label{thm:pbw}
Let $\cA$ carry an exhaustive filtration $F_\bullet$ with the properties:
$F$ is flat; the associated graded $\mathrm{gr}_F\cA$ is a classically
Koszul commutative chiral algebra; each finite conformal-weight window is
finite-dimensional and the resulting towers are Mittag--Leffler. Then
$q_{\cA}$ is a quasi-isomorphism: $\cA$ is chirally Koszul. In particular
the universal affine algebra $V_k(\fg)$ for every level $k$, the Virasoro
algebra $\Vir_c$ for every $c$, the universal principal
$\cW$-algebras, the free-field algebras and the lattice algebras are
chirally Koszul.
\end{theorem}

\begin{proof}
Filter the twisted tensor complex by PBW degree. Flatness identifies the
associated graded differential with the classical Koszul complex of
$\mathrm{gr}_F\cA$, which is acyclic by Priddy's theorem; the bracket and
level terms are of lower order. The spectral sequence collapses at $E_1$;
Mittag--Leffler gives strong convergence: the conformal-weight windows are
finite-dimensional, their tower has surjective transitions, and the
limit closes by Mittag--Leffler.
\end{proof}

Koszulness is a property of the universal algebra, not of its quotients:

\begin{theorem}\label{thm:non-koszul}
The Ising chiral algebra $M(3,4)$ lies outside the Koszul locus: the null-vector
quotient removes a bar-degree-one chain at weight $(p-1)(q-1)$ while its
image in degree two survives, producing $H^{2}(\barB)\neq0$ off the
diagonal. The simple admissible quotients $L_k(\mathfrak{sl}_3)$ at
denominator $q\geq3$ lie outside it as well: the Cartan subalgebra
contributes $\operatorname{rk}(\fg)$ permanent classes to $H^{2}(\barB)$.
Koszulness is logically independent of $C_2$-cofiniteness: $V_k(\fg)$ is
Koszul and generically not $C_2$-cofinite, while $C_2$-cofinite simple
quotients may fail Koszulness.
\end{theorem}

On the recognition locus the property has many faces. Six conditions are
equivalent through the single hub of PBW collapse: the existence of a
Koszul twisting datum; $q_{\cA}$ a quasi-isomorphism; the cobar comparison
$\Omega_X(\cA^{i})\to\cA$ a quasi-isomorphism; acyclicity of the two
twisted tensor complexes; collapse of the bar spectral sequence on the
diagonal; and diagonal concentration of
$\operatorname{Ext}_{\cA}(\mathbf{1},\mathbf{1})$. A seventh condition,
formality over the Swiss-cheese operad, is finer: it is equivalent to
the conjunction of the hub with membership in the Gaussian class of
Section~\ref{sec:shadow}, and the affine and Virasoro algebras, Koszul
and non-Gaussian, exhibit the gap between the six and the seventh.

\subsection{Curvature of the dual}

The dual of an algebra whose OPE has double poles is curved. For
$\mathcal{H}_k$ the dual is $(\operatorname{Sym}^{\mathrm{ch}}(V^{*}[1]),
m_0=-k\omega)$; the curvature $m_0$ is the level. Curved coalgebras are
objects of Positselski's coderived category of the second kind \cite{Pos}:
the derived category of comodules taken with respect to coacyclic rather
than acyclic objects, in which a curved bar--cobar comparison with
coacyclic cone is an equivalence and the comodule--contramodule
correspondence
$D^{\mathrm{co}}(C\text{-comod})\simeq D^{\mathrm{ctr}}(C\text{-contra})$
holds for any coalgebra with finite-dimensional graded pieces. Four
regimes organise the inversion theory, each with its own convergence
mechanism:
\begin{center}
\begin{tabular}{llll}
\emph{regime} & \emph{examples} & \emph{curvature} & \emph{mechanism}\\
\hline
quadratic & Heisenberg, lattice & $0$ & none needed\\
curved central & affine, Virasoro & $\kappa\,\omega_g$ &
coderived category\\
filtered complete & $\cW_N$ & central & adic completion\\
infinite type & $\cW_\infty$ & central & weight-completed towers
\end{tabular}
\end{center}
The distinction becomes structural at positive genus and is taken up in
Section~\ref{sec:genus}.

\subsection{Completed inversion}

The chain-level realisation of the bar--cobar equivalence holds in four
regimes, each with its own convergence mechanism, and the passage between
them is a single lemma applied four times.

\begin{theorem}[Completed inversion]\label{thm:completed}
Let $\cA$ carry a complete descending filtration whose finite quotients
$\cA_{\leq N}$ are finite-type, with surjective transition maps. Then the
completed bar and cobar functors exist with continuous differentials, and
the completed counit
$\widehat\Omega\,\widehat{\barB}(\cA)\to\cA$ is a quasi-isomorphism in
the square-zero total model: a bar word of total weight $w$ involves only
modes bounded by $w$, so each finite window stabilises; the Milnor exact
sequence
$0\to{\varprojlim}^{1}H^{m-1}\to H^{m}(\widehat{\phantom{x}})\to
\varprojlim H^{m}\to0$
then closes the limit, the $\varprojlim^1$ term dying on the surjective
tower. In the genuinely curved regime the same finite-stage argument
yields a coacyclic-cone comparison in the coderived category rather than
a quasi-isomorphism. The four regimes --- quadratic (Heisenberg,
lattice: no completion needed), curved-central (affine, Virasoro:
scalar curvature), filtered-complete ($\cW_N$: adic completion), and
infinite-type ($\cW_\infty$: weight-completed towers) --- are distinct,
and every completed statement in this paper declares its regime.
\end{theorem}

\subsection{Modules and the character identity}

On the quadratic genus-zero lane, Koszul duality extends to modules: the
categories of complete modules over $\cA$ and conilpotent comodules over
$\barB(\cA)$ are equivalent, with Ext exchanged; the numerical shadow of
the equivalence is the Koszul character identity
\begin{equation}\label{eq:character}
h_{\cA}(t)\;h_{\cA^{i}}(-t)\;=\;1 .
\end{equation}
For the rank-$d$ Heisenberg algebra with its Fock modules the exchange
is completely explicit: distinct Fock modules have no extensions, and
the bigraded Ext series $\prod_n(1+tq^{n})^{d}$ returns
$\eta(\tau)^{d}$ at $t=-1$ --- the character identity meeting the
Fredholm determinant of Section~\ref{sec:blocks}.

\section{The ordered theory and quantum groups}\label{sec:ordered}

\subsection{What averaging forgets}

The symmetric bar complex is the quotient of the ordered one by the action
of the symmetric groups. The quotient map is the Reynolds averaging
operator, and it is lossy in a way that can be computed exactly.

\begin{theorem}[Reynolds kernel]\label{thm:reynolds}
Let $V$ be a $d$-dimensional space of fields. The kernel of the averaging
map $\mathrm{av}_n\colon V^{\otimes n}\to\operatorname{Sym}^{n}V$ has
dimension
\begin{equation}
\dim\ker(\mathrm{av}_n)\;=\;d^{\,n}-\binom{n+d-1}{d-1},
\end{equation}
with generating function $\frac{1}{1-dt}-\frac{1}{(1-t)^{d}}$: the ordered
theory grows exponentially, its symmetric shadow polynomially.
\end{theorem}

\begin{proof}
Schur--Weyl decomposition of $V^{\otimes n}$ against the stars-and-bars
count of $\operatorname{Sym}^{n}V$; the Reynolds operator projects each
isotypic block onto its trivial part, and every non-trivial
$\Sigma_n$-isotypic lies in the kernel.
\end{proof}

The content of the kernel is graded by degree. At degree two the ordered
theory retains the binary collision residue
\begin{equation}\label{eq:rmatrix}
r_{\cA}(z)\;=\;\mathrm{Res}^{\mathrm{coll}}_{0,2}\bigl(\Theta^{E_1}_{\cA}\bigr),
\end{equation}
the classical $r$-matrix; averaging retains only its trace, the scalar
$\kappa$ of Section~\ref{sec:genus}, with a shift: for the affine algebra
the simple-pole self-contraction through the adjoint Casimir adds
$\dim\fg/2$, so that
$\kappa=\mathrm{av}(r)+\dim\fg/2=\dim\fg\,(k+\dvee)/2\dvee$. At degree
three the kernel contains the commutator $[\Omega_{12},\Omega_{23}]$ and
with it the Drinfeld associator: any splitting of the averaging map is a
choice of associator, and the set of splittings is a torsor over the
Grothendieck--Teichm\"uller group $\mathrm{GRT}_1(\bQ)$. At genus one a
further obstruction appears: the modular completion of the propagator
replaces $E_2$ by $E_2-3/(\pi\operatorname{Im}\tau)$, and the
non-holomorphic correction lives strictly beyond the holomorphic
theory; the ordered and
symmetric theories are bound together over the modular parameter,
the splitting obstruction being exactly this anomaly. Beyond degree
three the kernel carries the full pure-braid monodromy of the
Knizhnik--Zamolodchikov connection: individual conformal blocks as
vectors, of which the symmetric theory keeps only the count.

Quantum groups, in one sentence, are the $E_1$ memory of chiral algebras.

\subsection{The inversion principle}

Koszul duality acts on the ordered theory through Verdier duality on
ordered configuration spaces, and Verdier duality reverses the orientation
of collision cycles. Orientation reversal inverts the noncommutativity
datum.

\begin{theorem}[Inversion principle]\label{thm:inversion-principle}
On the ordered theory, Koszul duality inverts the braiding datum: for a
lattice algebra with $2$-cocycle $\varepsilon$ the dual carries
$\varepsilon^{-1}$; for a braided input with $R$-matrix $R(z)$ the dual
carries $R(z)^{-1}$; on quantum parameters, $q\mapsto q^{-1}$. The
operadic form of the statement is
$(\mathrm{Ass}^{\mathrm{ch}})^{!}\cong\mathrm{Ass}^{\mathrm{ch}}
\otimes\mathrm{sgn}$.
\end{theorem}

\begin{proof}
Verdier duality exchanges the $j_{*}$-extension with logarithmic
forms and the $j_{!}$-extension with boundary currents, and on
ordered configuration spaces it reverses the co-orientation of each
collision cycle. The pairing of a cycle with the noncommutativity
datum --- cocycle value, $R$-matrix, braiding phase --- is therefore
inverted, and the sign representation records the reversal at the
operadic level.
\end{proof}

This has a consequence worth isolating. The level inversion
$k\mapsto-k-2\dvee$ of the symmetric theory (the unique affine involution
fixing the critical level and compatible with $\kappa$) and the braiding
inversion $q\mapsto q^{-1}$ of the ordered theory are the same operation
seen at two depths, since $q=e^{\pi i/(k+\dvee)}$.

\subsection{The Yangian}

The Yangian enters not as an example but as the value of the ordered
duality functor. In the RTT presentation, with Yang matrix
$R_{\hbar}(u)=1-\hbar P/u$, the defining relation
$R_{12}(u{-}v)T_1(u)T_2(v)=T_2(v)T_1(u)R_{12}(u{-}v)$ equips the completed
Yangian with a quadratic $E_1$-factorization structure whose braiding is
genuinely non-symmetric: its bar complex exists only on ordered
configurations.

\begin{theorem}\label{thm:yangian}
For the Yang kernel in type $A$, the orthogonal complement of the
RTT relation space under the trace pairing is the relation space of the
inverse kernel, and
$R_{\hbar}(u)^{-1}=(1-\hbar^{2}/u^{2})^{-1}R_{-\hbar}(u)$. Each finite
RTT window therefore satisfies
$A_N(R_{\hbar})^{!}\cong A_N(R_{-\hbar})$, and since a bar word of total
weight $w$ involves only modes $\leq w$, the windows stabilise and the
completed statement follows:
\begin{equation}
Y_{\hbar}(\mathfrak{sl}_2)^{!}\;\cong\;Y_{-\hbar}(\mathfrak{sl}_2),
\qquad \hbar=\frac{1}{k+\dvee}.
\end{equation}
\end{theorem}

\begin{proof}
The RTT relation presents the window as the image of the operator
$X\mapsto R_{12}X-\sigma(X)R_{12}$; a direct computation in the
fundamental square identifies the orthogonal complement of this image
under the trace pairing with the corresponding image for
$R^{-1}$, and unitarity
$R_{12}(u)R_{21}(-u)=(\hbar^{2}-u^{2})\,\mathrm{id}$ converts
$R^{-1}$ into $R_{-\hbar}$ up to the central scalar, which preserves
the relation span. Quadratic duality on each window is therefore the
sign flip; the windows stabilise as stated, and the Milnor sequence
completes the limit.
\end{proof}

Three remarks fix the scope. First, the identification of the ordered
Koszul dual of $\widehat{\fg}_k$ itself with the Yangian requires, beyond
Theorem~\ref{thm:yangian}, a strictification of the quasi-associative
structure; the obstruction lives in root spaces and vanishes when all root
multiplicities are one --- the mechanism fails precisely for Kac--Moody
algebras with imaginary roots, which is where the Borcherds algebras of
Section~\ref{sec:arithmetic} live. Second, the ordered double dual loses
the level: $(\cA^{!})^{!}\simeq U(\fg[t])$, not $\widehat{\fg}_k$. The
central extension is held by the closed (holomorphic) colour through the
double pole, the Yangian datum by the open (topological) colour through
the simple pole; only the two-coloured double duality is involutive. The
two colours of the Swiss-cheese operad carry disjoint information,
each indispensable. Third, at genus one the $B$-cycle
monodromy of the elliptic propagator exponentiates
$\hbar$ to $q=e^{\pi i/(k+\dvee)}$; this is a monodromy statement about
local systems; the algebra isomorphism with $U_q(\fg)$ is a separate
statement, and monodromy is what the geometry supplies.

\subsection{Monodromy, Langlands duality, and the evaluation core}

The degree-three content of the ordered theory is monodromy. The shadow
connection of the affine algebra is the Knizhnik--Zamolodchikov
connection; its monodromy generates the braid-group representations of
the quantum group at $q_{\mathrm{KL}}=e^{\pi i/(k+\dvee)}$, and the
level reflection acts as $q\mapsto q^{-1}$, which is
Theorem~\ref{thm:inversion-principle} again. The identification of the
quantum parameter is overdetermined: for $\widehat{\mathfrak{sl}}_2$ the
braiding eigenvalues on the two channels of the fundamental square are
$-q^{-3/2}$ and $q^{1/2}$, and four independent computations --- the
KZ monodromy, the evaluation-module braiding at
$u=\cot\frac{\pi}{k+2}$, the fusion-category twist, and the
$B$-cycle monodromy at genus one --- return the same $q$; the number
$1/(k+\dvee)$ is the unique coupling compatible with the Sugawara
construction, and every road leads to it. Three statements calibrate
what is known beyond type $A$.

First, the categorical Clebsch--Gordan input holds in complete
generality: for every simple $\fg$, the tensor product of a fundamental
evaluation module with a negative prefundamental in the shifted
category \cite{HJ} $\mathcal{O}$ splits according to the $q$-character,
$V_{\omega_i}(a)\otimes L^{-}_{i}(b)\cong\bigoplus_{\mu}
L^{-}_{i}(\text{shift}=\mu)$, at generic parameters: the summand
$\ell$-weights are multiplicity-free and pairwise unlinked, so the
filtration by singular vectors splits. On the category generated by
evaluation modules, ordered Koszul duality is therefore established for
all simple types; its extension from the evaluation core to the full
category is a completion problem, open beyond type $A$, and in type $A$
reduced to a single compact-generation statement on the Baxter locus
$b=a-\tfrac{\lambda+1}{2}$.

Second, the Langlands form of the duality is exactly as strong as its
inputs: granting the lacing isomorphism of Frenkel--Hernandez \cite{FH} and a
Chevalley lift of $\hbar\mapsto-\hbar$, the inverse-kernel duality of
Theorem~\ref{thm:yangian} transports to
$Y_{\hbar}(\fg)^{!}\simeq Y_{-\hbar\,r_\fg}(\fg^{L})$, with
$r_\fg$ the lacing number: self-dual for the simply-laced types, and
exchanging $B_r$ with $C_r$. For the exceptional types the duality is
a finite construction in seven items --- kernel, projectors,
Poincar\'e--Birkhoff--Witt basis, relation-space isomorphism, scalar
gauge, tower compatibility, chiralisation --- whose completion yields
$\widehat Y_{\hbar}(\fg)^{!}\cong\widehat Y_{-\hbar}(\fg)$. The
complementarity equation
$\kappa(V_k(\fg))+\kappa(V_{k^{\vee}}(\fg^{\vee}))=0$ has the unique
solution $k^{\vee}=-k-2\dvee$: the level of the dual theory is fixed
by the vanishing of the total curvature.

Third, the super form: for $Y_{\hbar}(\mathfrak{sl}(m|n))$ with
$m\neq n$ the supertrace complementarity
$\kappa^{\mathrm{str}}+\kappa^{\mathrm{str},!}=0$ holds under
$k\mapsto-k-2(m-n)$; the Berezinian-normalised sum is
$\max(m,n)$: a conjecture, supported by the closed forms at low rank.

\subsection{The Etingof--Kazhdan triangle}

On the ordered Koszul locus, at a fixed associator, three structures
on the bar coalgebra determine one another: a vertex $R$-matrix
$S(z)$ satisfying the Etingof--Kazhdan axioms; the curved
$A_\infty$-structure of the ordered bar complex; and a chiral
coproduct, coassociative up to the associator, on the double. The
equivalence triangle is verified in full for the
$\mathfrak{sl}_2$ Yangian and the affine algebras, and the
quantization functor of Etingof--Kazhdan \cite{EK} is, in this
language, a choice of splitting of the averaging map --- which is why
its ambiguity is the same Grothendieck--Teichm\"uller torsor. The
triangle is the precise sense in which a chiral algebra with a
braiding \emph{is} a quantum group: three definitions, one object,
one torsor of identifications.

\subsection{Two normalisations}

The collision residue of the affine algebra is $r=k\Omega_{\mathrm{tr}}/z$
in the trace form; the Knizhnik--Zamolodchikov connection carries
$\Omega/\bigl((k+\dvee)z\bigr)$ with $\Omega=2\dvee\,\Omega_{\mathrm{tr}}$.
Their ratio is $k(k+\dvee)/2\dvee$: no rescaling independent of the level
identifies them, as evaluation at $k=0$ shows. The first vanishes at
$k=0$ and is finite at the critical level; the second is the reverse.
Every formula in this paper declares which of the two it uses.

\subsection{The dictionary}

The division of labour between the two theories is exact:
\begin{center}
\begin{tabular}{ll}
\emph{the symmetric theory keeps} & \emph{the ordered theory keeps}\\
\hline
the modular characteristic $\kappa$ & the $r$-matrix $r(z)$\\
the shadow tower $S_r$ & the associator $\Phi$\\
the genus expansion $F_g$ & the braid monodromy\\
characters and Verlinde numbers & conformal blocks as vectors\\
the conductor $\varkappa$ & the quantum group
\end{tabular}
\end{center}
Every left-hand entry is the average of its right-hand neighbour, and
each row of the left column is computed in this paper from the row on
the right.

\section{Higher genus: curvature and the Maurer--Cartan element}
\label{sec:genus}

\subsection{The prime form and the curved differential}

On a curve of genus $g\geq1$ the prime form replaces the difference
$z_1-z_2$: it is its $E(z_1,z_2)$, a section of
$K^{-1/2}\boxtimes K^{-1/2}$ vanishing simply on the diagonal, and the
propagator becomes
\begin{equation}\label{eq:genus-propagator}
\eta^{(g)}_{12}
= d\log E(z_1,z_2)
+\pi\sum_{j,k=1}^{g}(\operatorname{Im}\Omega)^{-1}_{jk}
\operatorname{Im}\!\Bigl(\int_{z_2}^{z_1}\!\omega_j\Bigr)\,\omega_k(z_1),
\end{equation}
where $\Omega$ is the period matrix. The correction term is forced: the
prime form is quasi-periodic around $B$-cycles, and \eqref{eq:genus-propagator}
is the unique completion of $d\log E$ that is single-valued. Single-valuedness
exchanges holomorphy for curvature.

\begin{theorem}[Curvature]\label{thm:curvature}
Let $\cA$ be a chiral algebra of the standard landscape and let
$\Dfib$ be the bar differential built from \eqref{eq:genus-propagator} on
a family of genus-$g$ curves. Then $\Dfib$ is a curved differential:
$\Dfib^{2}=[m_0^{(g)},-]$ for a curvature element $m_0^{(g)}$ whose
scalar diagonal is
\begin{equation}
\operatorname{tr}_{\mathrm{diag}}\bigl(m_0^{(g)}\bigr)
=\kappa(\cA)\cdot\omega_g,
\qquad \omega_g=c_1(\bE_g)
\end{equation}
the first Chern class of the Hodge bundle. The scalar $\kappa(\cA)$ is
computed at genus one by two channels: the $B$-cycle monodromy defect
$-2\pi i$ of the propagator acting through the double pole, and a one-loop
self-contraction through the simple pole with symmetry factor $\tfrac12$
and adjoint Casimir $2\dvee$. The resulting values are
\begin{equation}\label{eq:kappa-census}
\begin{aligned}
\kappa(\mathcal{H}_k)&=k, &
\kappa\bigl(V_k(\fg)\bigr)&=\frac{(k+\dvee)\dim\fg}{2\dvee}, &
\kappa(\Vir_c)&=\tfrac{c}{2},\\
\kappa(\cW_N)&=c\,(H_N-1), &
\kappa(\beta\gamma_\lambda)&=6\lambda^{2}-6\lambda+1, &
\kappa(V_\Lambda)&=\operatorname{rk}\Lambda,
\end{aligned}
\end{equation}
with $H_N=\sum_{j\leq N}1/j$, and $\kappa(bc_\lambda)=-\kappa(\beta\gamma_\lambda)$;
$\kappa$ is additive under tensor product. At the critical level
$k=-\dvee$ the two channels cancel: $\kappa=0$.
\end{theorem}

\begin{proof}[Proof at genus one]
Two computations. The corrected propagator on the torus is
$\bigl[\theta_1'/\theta_1(z)+2\pi i\operatorname{Im}z/
\operatorname{Im}\tau\bigr]dz$; it is $1$-periodic and acquires the
$B$-cycle defect $-2\pi i$. Squaring the differential, the double-pole
channel transports the level around the $B$-cycle and produces
$k\,\dim\fg\cdot\omega_1$ from the defect; the simple-pole channel
closes a self-contraction whose Casimir is the Killing form,
$\sum f^{ac}_{\;\;d}f^{bc}_{\;\;d}=2\dvee\kappa^{ab}$, with symmetry
factor $\tfrac12$, producing $\dvee\dim\fg\cdot\omega_1$. Normalising
by $2\dvee$ gives $\kappa=(k+\dvee)\dim\fg/2\dvee$. The identity
$\partial_\tau\log\eta=\tfrac{\pi i}{12}E_2$ identifies the scalar
class with the quasi-modular anomaly of $E_2$, which is the analytic
form of the statement that the curvature is intrinsically
quasi-modular.
The remaining rows of \eqref{eq:kappa-census} are the same two-channel
computation with the appropriate Casimir data; additivity is manifest
channel by channel.
\end{proof}

The two channels of the proof define a canonical decomposition
$\kappa=\kappa_{\mathrm{dp}}+\tfrac12\dim\fg$: the double-pole part
is the averaged $r$-matrix, the simple-pole part the Sugawara shift,
and the level reflection $k\mapsto-k-2\dvee$ acts on the pair as the
shear fixing the shift and negating the total --- $\kappa$ is the
$(-1)$-eigenvector of the reflection, which is how the dual level in
Theorem~\ref{thm:complementarity} is found. Two refinements of
Theorem~\ref{thm:curvature} deserve record. At genus one
the double-pole channel alone contributes, since holomorphic differentials
are constant; from genus two onward $\partial_w\omega_j\neq0$ and the level
enters the curvature directly. And the simply-laced level-one algebras
collapse: their root currents are composites of $\operatorname{rk}\fg$
Heisenberg fields, the simple-pole channel dies, and
$\kappa(\widehat\fg_1)=\operatorname{rk}\fg$ --- the unique crossing
between classes in the classification of Section~\ref{sec:shadow}.

\begin{theorem}[Fay's identity as the higher-genus Borcherds identity]
\label{thm:fay}
On a smooth curve of genus $g$, the vanishing of the square of the
period-corrected differential $D_g$ reduces, on triple-collision
strata, to the trisecant identity of Fay for the prime form and the
Szeg\H{o} kernel: writing $\eta^{(g)}=h+a$ for the meromorphic and
real-analytic parts of \eqref{eq:genus-propagator}, the cyclic sum
$\sum_{\mathrm{cyc}}h_{12}h_{23}$ is killed by the differentiated
trisecant identity, the pure $a\!\cdot\!a$ terms vanish as
$(2,0)$-forms on a curve, and the cross terms assemble the Bergman
kernel, whose diagonal residue is the curvature of
Theorem~\ref{thm:curvature}. The Borcherds identity governs the
algebra at every genus; Fay's identity governs the geometry, and the
two meet only through the scalar $\kappa$.
\end{theorem}

\begin{proof}
Squaring $D_g$ and sorting by type: the meromorphic square is the
cyclic sum $\sum_{\mathrm{cyc}}h_{12}h_{23}$, which the
differentiated trisecant identity of Fay \cite{Fay} expresses through
even Weierstrass data and antisymmetric form factors, hence kills; the
$a\!\cdot\!a$ terms are $(2,0)$-forms on a curve and vanish; and
$\bar\partial a$ is the Bergman kernel, so the cross terms reduce to
the diagonal residue evaluated in Theorem~\ref{thm:curvature}.
\end{proof}

\begin{theorem}[Locality]\label{thm:locality}
Collision residues see only the logarithmic kernel: the residue of the
second term of \eqref{eq:genus-propagator} along any diagonal vanishes,
so the bar differential restricted to any fixed smooth curve equals its
genus-zero form. Every algebraic invariant of the bar complex ---
Koszulness, the obstruction tower, the classification --- is genus-zero
data. The genuinely modular content of a chiral algebra is concentrated
in the single object constructed next.
\end{theorem}

\subsection{The Maurer--Cartan element}

Flatness is restored by geometry. The Gauss--Manin connection of the
family cancels the curvature, $D_g=\Dfib+\nabla^{\mathrm{GM}}$ satisfies
$D_g^{2}=0$, and the bar complexes over all genera assemble into an
algebra over the Feynman transform of the modular operad
$\{\Mbar_{g,n}\}$ of Getzler and Kapranov \cite{GK}.

\begin{definition}
The \emph{modular convolution Lie algebra} of $\cA$ is
\begin{equation}
\gmod_{\cA}
=\prod_{2g-2+n>0}
\operatorname{Hom}_{\Sigma_n}\!\bigl(C_{*}(\Mbar_{g,n}),\,
\operatorname{End}_{\cA}(n)\bigr),
\end{equation}
with differential induced by the boundary of $\Mbar_{g,n}$ and bracket by
one-edge sewing; the weight filtration by $2g-2+n$ makes it pronilpotent.
\end{definition}

\begin{theorem}[The master object]\label{thm:mc}
Let $D_{\cA}$ be the total genus-completed bar differential and $d_0$ its
tree-level part. Then $\Theta_{\cA}:=D_{\cA}-d_0\in\gmod_{\cA}$ satisfies
\begin{equation}\label{eq:MC}
D\,\Theta_{\cA}+\tfrac12\,[\Theta_{\cA},\Theta_{\cA}]=0
\end{equation}
identically: the equation is the expansion of $D_{\cA}^{2}=0$, which holds
because the boundary of the boundary of $\Mbar_{g,n}$ is empty. No
equation holds term by term, and convergence
of $\Theta_{\cA}$ in the weight completion is automatic.
\end{theorem}

Theorem~\ref{thm:mc} is the central statement of the paper, and its proof
is its point: the Maurer--Cartan equation of two-dimensional chiral
physics is the transcription of $\partial^{2}=0$ on Deligne--Mumford
space. Everything constructed below --- the obstruction tower, the
complementarity law, the genus expansion, the $r$-matrix, the
Knizhnik--Zamolodchikov connection --- is a projection of
$\Theta_{\cA}$ along a filtration. The Maurer--Cartan equation at
$(g,n)$, expanded, reads: internal differential, plus separating
clutching (the bracket), plus non-separating self-sewing (a
Batalin--Vilkovisky operator $\hbar\Delta$), equals zero; it determines
$\Theta^{(g,n)}$ from $(0,3)$ and $(0,4)$ by induction, uniquely up to
gauge on the Koszul locus. The induction replaces the sum over
topologies of three-dimensional gravity; there is no divergent sum to
regularise.

For the affine algebras the element is strict in closed form: on the
one-channel lane $\Theta^{\mathrm{str}}=\kappa\cdot\mu\otimes\Lambda$,
with $\mu$ the invariant functional and
$\Lambda=\sum_{g\geq1}(-1)^{g}\lambda_g$ the total Hodge class ---
closed because the Killing three-cocycle is closed, with vanishing
bracket by the Jacobi identity, exact at every genus.
The simplest gravitational background of the theory is the level
times the total Chern class of the Hodge bundle, in closed form.

\subsection{The two levels of the master equation}

The proof of Theorem~\ref{thm:mc} lives at two levels, and separating
them is part of the statement. At the convolution level, $D^{2}=0$ is
pure topology: the differential is composition against the singular
boundary of $\Mbar_{g,n}$, and its square vanishes because
$\partial^{2}=0$; no geometry of the fibres enters. At the ambient
level --- where the differential carries, besides clutching, the
collision faces, the logarithmic degenerations, and the planted-forest
corrections of the relative compactification --- the square-zero
identity is the signed residue--pushforward calculus of the
logarithmic Fulton--MacPherson geometry, and it holds granted that
calculus; the fixed-smooth-curve core stands on its own. Every scalar
consequence drawn in this paper factors through the convolution
level.

Expanded, the master equation is the quantum master equation of
Batalin and Vilkovisky. The separating boundary is the bracket; the
non-separating boundary is a second-order operator $\Delta$ raising
genus and lowering arity by two; and \eqref{eq:MC} reads
$\hbar\Delta S+\tfrac12\{S,S\}=0$ for $S=\Theta_{\cA}$. The clutching
maps factorise the element,
$\xi_{h}^{*}\Theta^{(g)}=\sum_{g_1+g_2=g}
\Theta^{(g_1)}\star\Theta^{(g_2)}$ and
$\xi_{\mathrm{irr}}^{*}\Theta^{(g+1)}=\Delta\,\Theta^{(g)}$, and
Verdier duality intertwines the elements of an algebra and its dual,
$\mathbb{D}(\Theta_{\cA})=\Theta_{\cA^{!}}$ --- an intertwining, not a
negation: $\kappa+\kappa^{!}=0$ holds on the affine and free lanes and
fails already for Virasoro, where the sum is the conductor of
Section~\ref{sec:conductor}.

\subsection{The primitive kernel}

The Maurer--Cartan element compresses to a primitive kernel: the sum
$\mathfrak{K}$ of corolla terms, one for each stable topological
type, satisfying the primitive master equation
\begin{equation}
d\mathfrak{K}+\mathfrak{K}\star\mathfrak{K}
+\hbar\,\Delta\mathfrak{K}=0,
\end{equation}
with $\star$ the one-edge sewing and $\Delta$ the self-sewing. The
first values are structural: the modular characteristic is the trace
of the self-sewn binary corolla,
$\kappa(\cA)=\operatorname{tr}_{\mathrm{cyc}}\Delta\,
\mathfrak{K}_{0,2}$ --- the genus-one tadpole of the two-point
function --- and the genus-one one-point kernel is
$\mathfrak{K}_{1,1}=-P\circ\Delta\,\mathfrak{K}_{0,3}$, the
propagator applied to the self-sewn three-point corolla. Genus by
genus, the master equation writes each kernel as a polynomial in
lower ones: the entire modular structure of a chiral algebra is the
free resolution of its operator product expansion by stable curves.

\subsection{Five points, sewn}

The first genuinely global test of the master equation is the
five-point sphere, and it passes at chain level. On the formal
completion of $\Mbar_{0,5}$, with its ten boundary divisors, the
five-point correlator of a Koszul algebra extends uniquely across
every divisor by the sewing rule
\begin{equation}
\Phi_5\big|_{D_{S,S^{c}}}
=\sum_{a,b}\Phi_3(\xi_S;e_a)\,\eta^{ab}\,\Phi_4(e_b;\xi_{S^{c}}),
\end{equation}
$\eta^{ab}$ the inverse Shapovalov form, one rule serving all ten
divisors, which the $\Sigma_5$-action permutes transitively; the
arity-five bar
differential is the pullback of the Arnold connection, its square
vanishing by the three-term relation together with the infinitesimal
braid relations; and the extension is compatible with the genus-one
clutching charts through the elliptic trace. The pentagon of
$\Mbar_{0,5}$ is where associativity of sewing is first a theorem
rather than a definition, and the computation here is the base case
of the induction of Theorem~\ref{thm:mc}.

\subsection{The coderived ambient}

At positive genus the completion becomes part of the structure, and
the following theorem locates it exactly.

\begin{theorem}\label{thm:coderived}
For the Virasoro and $\cW_N$ algebras the identification of the
Batalin--Vilkovisky complex with the bar complex is false at the chain
level in the category of bounded direct-sum complexes: the harmonic
components of the propagator produce discrepancies
$\delta_r=c_r(\cA)\,m_0^{\lfloor r/2\rfloor-1}$, $r\geq4$, with
$c_4(\Vir_c)=10/[c(5c+22)]\neq0$, and no contracting homotopy absorbs
them. The discrepancies are multiples of powers of the curvature; their
cone is coacyclic. Consequently the identification holds in
Positselski's coderived category, equivalently on the pro-object,
adic, or weight-completed presentations, which are canonically
equivalent; the element
$\xi=(e_2,e_3,\dots)$ with $e_k=s^{-1}L_{-k}\mathbf{1}\otimes
s^{-1}L_{-k}\mathbf{1}$ lies in the product completion and not in the
direct sum: the completion is intrinsic.
\end{theorem}

\begin{proof}
Hodge-decompose the genus-$g$ propagator into the logarithmic kernel,
an exact part, and a harmonic part. The exact part is a chain
homotopy, so the whole discrepancy between the two differentials
concentrates in the harmonic directions, where degree counting on
the Siegel space forces each correction to be a multiple of a power
of the central curvature. A cone built from curvature-torsion is
resolved by the one-variable Koszul complex of $m_0$, hence
coacyclic; inverting coacyclics --- the definition of the coderived
category --- makes the comparison an equivalence. The three completed
presentations are identified by the finite-window stabilisation of
Theorem~\ref{thm:completed}.
\end{proof}

Theorems~\ref{thm:curvature}--\ref{thm:coderived} say one thing:
nilpotence is the logarithm of periodicity. The bar construction is the
categorical logarithm --- it takes tensor products to direct sums --- the
cobar construction is the exponential, the Koszul locus is the radius of
convergence, $\kappa$ is the infinitesimal generator of monodromy, and
the coderived category is the analytic continuation.

\section{The shadow tower and the classification}\label{sec:shadow}

\subsection{The tower}

Restrict $\Theta_{\cA}$ to a one-dimensional primary line $L$ and expand:
$\Theta_{\cA}|_{L}=\sum_{r\geq2}S_r\,x^{r}$. The coefficients $S_r$ are the
\emph{shadow tower} of $\cA$; they are the complete set of scalar
obstructions to extending the truncations
$\Theta^{\leq2}\to\Theta^{\leq3}\to\cdots$, and by
Theorem~\ref{thm:locality} they are genus-zero data. The first three have
independent names and meanings: $S_2=\kappa$ is the modular characteristic
of Theorem~\ref{thm:curvature}; $S_3$ is the cubic collision residue; $S_4$
is the quartic contact term, computed for the Virasoro algebra by the
Zamolodchikov norm: the unique quasi-primary at level four is
$\Lambda=\,:TT:-\tfrac{3}{10}\partial^{2}T$ with
$\langle\Lambda|\Lambda\rangle=c(5c+22)/10$, and
\begin{equation}\label{eq:S4}
S_3(\Vir_c)=2\quad(\text{independent of }c),
\qquad
S_4(\Vir_c)=\frac{10}{c\,(5c+22)}.
\end{equation}

The quartic seed is a Gram computation. At level four the Verma module
of $\Vir_c$ has the two-dimensional basis $\{L_{-2}^{2}v,\,L_{-4}v\}$
modulo derivatives, with Gram matrix
$\begin{pmatrix}5c&3c\\3c&c(c+8)/2\end{pmatrix}$ of determinant
$c^{2}(5c+22)/2$; the quasi-primary direction is $\Lambda$, its norm is
the Zamolodchikov value, and $S_4$ is its inverse. The two zeros of the
Gram determinant, $c=0$ and the Lee--Yang point $c=-22/5$, are the only
poles of every shadow coefficient of the family: the Kac determinant at
level four governs the arithmetic of the entire tower. The quintic
obstruction is the first genuinely nonlinear one,
$o^{(5)}=\{S_3,S_4\}_{H}=480/[c^{2}(5c+22)]\neq0$, and its
non-vanishing is the one-line proof that the Virasoro tower is
infinite.

\begin{theorem}[Algebraicity]\label{thm:riccati}
Set $H(t)=\sum_{r\geq2}rS_rt^{r}$, $\alpha=S_3$, $\Delta=8\kappa S_4$. The
Maurer--Cartan recursion on the line $L$ is equivalent to the single
algebraic identity
\begin{equation}\label{eq:riccati}
H(t)^{2}\;=\;t^{4}\,Q(t),
\qquad
Q(t)\;=\;(2\kappa+3\alpha t)^{2}+2\Delta\,t^{2}.
\end{equation}
Three numbers $(\kappa,\alpha,S_4)$ determine the entire tower:
$S_r=\tfrac1r\,[t^{\,r-2}]\sqrt{Q(t)}$. The coefficients are rational in
the parameters of the family; for the Virasoro algebra the denominators
are confined to $c^{\,r-3}(5c+22)^{\lfloor(r-2)/2\rfloor}$, and
$S_5=-48/[c^{2}(5c+22)]$.
\end{theorem}

\begin{proof}
Write $a_n=(n+2)S_{n+2}$. The degree-$r$ component of \eqref{eq:MC} on $L$
is the vanishing of the convolution coefficient
$\sum_{i+j=m}a_ia_j$ for $m\geq3$; isolating the terms with $i=0$ or
$j=0$ gives $4r\kappa S_r$ against the restricted bracket sum, which is
the Taylor recursion for $\sqrt{Q}$; the coefficients of orders
$0,1,2$ are the definition of $Q$.
\end{proof}

\begin{theorem}[Depth gap]\label{thm:depth-gap}
Let $r_{\max}$ be the largest $r$ with $S_r\neq0$. On a single primary
line, $r_{\max}\in\{2,3,\infty\}$: the tower terminates if and only if
$Q$ is a perfect square, and a square-root of a non-square quadratic never
truncates. Depth four is realised only by a family phenomenon: on the
$\beta\gamma_\lambda$ family, $S_3=0$ and $S_4=-5/12$ with $S_{r\geq5}=0$
by rank-one rigidity of the contact stratum, a two-parameter statement
invisible on any single line. The list of depths is complete.
\end{theorem}

The discriminant $\Delta=8\kappa S_4$ decides everything: $\Delta=0$
forces termination, $\Delta\neq0$ forces the infinite tower.

\subsection{The shadow connection}

The tower carries its own differential equation. The logarithmic
connection $\nabla^{\mathrm{sh}}=d-\frac{Q'(t)}{2Q(t)}\,dt$ is flat
with residue $\tfrac12$ at each simple zero of $Q$ and monodromy $-1$
around it --- the Koszul sign realised as monodromy --- and the
generating function satisfies the Picard--Fuchs equation
$2Qf''+Q'f'-Q''f=0$ of the genus-zero spectral curve $y^{2}=Q(t)$.
The tower is therefore a family of periods of a rational curve with
four branch points, its large-order behaviour a WKB analysis of the
associated Schr\"odinger operator, and its classical Voros period is
$\pi$, universally: the semiclassical structure of the obstruction
theory is the same for every chiral algebra, and only the position of
the branch points --- the three seeds --- distinguishes one family
from another.

\subsection{The classification}

\begin{theorem}[Classification]\label{thm:classification}
The chiral algebras of the standard landscape fall into four classes and
one critical stratum, according to the shadow tower, and membership is
forced by an identifiable mechanism:
\begin{center}
\begin{tabular}{lcccl}
\emph{Class} & $r_{\max}$ & $S_3$ & $S_4$ & \emph{mechanism and archetypes}\\
\hline
$\mathsf{G}$ & $2$ & $0$ & $0$ &
pure double poles: Heisenberg, fermion, lattice\\
$\mathsf{L}$ & $3$ & $\neq0$ & $0$ &
Jacobi identity kills $S_4$: affine $\widehat{\fg}_k$, $k\neq-\dvee$\\
$\mathsf{C}$ & $4$ & $0$ & $\neq0$ &
quartic contact, family stratum: $\beta\gamma$, $bc$\\
$\mathsf{M}$ & $\infty$ & $\neq0$ & $\neq0$ &
$T$ is a strong generator: $\Vir$, $\cW_N$, $V^{\natural}$\\
$\mathsf{FF}$ & --- & --- & --- &
critical level $k=-\dvee$: $\kappa=0$, uncurved, non-Koszul
\end{tabular}
\end{center}
For class $\mathsf{L}$ the vanishing of $S_4$ is the Jacobi identity in
the form
\[
\textstyle\sum_{(ij|kl)}\kappa_\fg([x_i,x_j],[x_k,x_l])
=\langle\mathrm{Jac},x_4\rangle=0,
\]
uniform over all simple types, the dual Coxeter number carrying the
non-simply-laced cases. Class $\mathsf{M}$ is characterised by
self-reference: $r_{\max}=\infty$ if and only if $T$ appears in its own
first product, $T\in T_{(1)}T$. The critical stratum is genuinely
distinct: at $k=-\dvee$ the bar complex is uncurved and yet not Koszul,
the Feigin--Frenkel centre
$\operatorname{Fun}\bigl(\mathrm{Op}_{{}^{L}\fg}(D)\bigr)$ spreading its
cohomology off the diagonal \cite{FF}. Uncurved and Koszul are
independent conditions.
\end{theorem}

\begin{proof}[Proof of the class assignments]
Class $\mathsf{G}$: a pure double-pole algebra has scalar residue
differential; all products of shadows beyond degree two vanish
identically. Class $\mathsf{L}$: the cubic shadow is the structure
constant tensor, and the quartic obstruction is the Hessian
$\tfrac12\{S_3,S_3\}$, whose value on any test quadruple is a sum of
Killing-paired double brackets; the Jacobi identity is exactly its
vanishing, as the displayed $\mathfrak{sl}_2$ computation
$\tfrac{2}{k+2}+0-\tfrac{2}{k+2}=0$ shows in coordinates. Class
$\mathsf{C}$: on the $\beta\gamma_\lambda$ family the cubic vanishes
by ghost-number parity and the quintic and beyond vanish by rank-one
rigidity --- the deformation direction is one-dimensional, and every
higher transferred product factors through the vanishing bracket of
degree two. Class $\mathsf{M}$: with $T\in T_{(1)}T$ the seeds
$\kappa$, $S_3$, $S_4$ are all nonzero, $Q$ has nonzero discriminant,
and Theorem~\ref{thm:depth-gap} forces the infinite tower; the
quintic value $480/[c^{2}(5c+22)]$ exhibits the first nonvanishing
obstruction explicitly. The critical stratum is
Example~\ref{ex:km} at $k=-\dvee$ together with the Feigin--Frenkel
description of the centre.
\end{proof}

\begin{theorem}[The bell]\label{thm:bell}
Three integers attached to $\cA$ coincide: the shadow depth $r_{\max}$;
the $A_\infty$-depth of the minimal model of $\cA$ (the largest $n$ with
$m_n\neq0$); and the $L_\infty$-formality level of the genus-zero
convolution algebra. A chiral algebra's homotopical complexity is a
single number, and class $\mathsf{G}$ is precisely the formal class.
\end{theorem}

\begin{proof}
Homotopy transfer identifies the spectral-sequence differentials
$d_r$ with the transferred operations $m_{r+1}$, so the first
nonvanishing higher product and the depth of the obstruction tower
coincide; the Kadeishvili tree formula and the stable-graph expansion
of the tower run over the same trees with the same propagators, giving
the third equality.
\end{proof}

Three separations complete the picture. The classification is a
classification of towers, not of central charges: the Moonshine module and
the Leech lattice algebra share $c=24$ and separate as
$(\kappa,\text{class})=(12,\mathsf{M})$ versus $(24,\mathsf{G})$. The
class and conductor are independent refinements
(Section~\ref{sec:conductor}):
Virasoro and $\cW_3$ are both class $\mathsf{M}$ with conductors $13$ and
$250/3$. And three pole-invariants are independent: the maximal generator
pole $p_{\max}$, its $d\log$-absorption $k_{\max}=p_{\max}-1$, and
$r_{\max}$ --- maximally split by $\beta\gamma$, where they equal
$(1,0,4)$. Drinfeld--Sokolov reduction is the class escalator
$\mathsf{L}\to\mathsf{M}$: gauge symmetry is turned into gravity by
quantised Hamiltonian reduction, and this is a construction, not a
metaphor.

\subsection{Arithmetic of the Virasoro tower}

The Virasoro tower rewards exact computation. Beyond \eqref{eq:S4} and
$S_5=-48/[c^{2}(5c+22)]$, the denominators obey the confinement law
$c^{\,r-3}(5c+22)^{\lfloor(r-2)/2\rfloor}$: every pole of every shadow
coefficient sits at $c=0$ or at the Lee--Yang point $c=-22/5$, the two
zeros of the level-four Gram determinant $c^{2}(5c+22)/2$ --- the Kac
determinant governing the arithmetic of the entire tower from its first
nontrivial case. The dual towers satisfy the reciprocity
\begin{equation}
\Delta(c)+\Delta(26-c)\;=\;\frac{6960}{(5c+22)(152-5c)},
\qquad 6960=2^{4}\cdot3\cdot5\cdot29,
\end{equation}
whose two poles are the Lee--Yang points of the dual pair. The branch
field of the tower is $\bQ\bigl(\sqrt{-5(5c+22)}\bigr)$, preserved by
the Koszul involution only at $c=1$ and the self-dual point $c=13$.
Finally, the genus at which a given coefficient first contributes to
the free energy is $g_{\min}(S_r)=\lfloor r/2\rfloor+1$: the tower is
graded by visibility, and the quartic term is a genus-three phenomenon
on the nose.

The closed forms continue: with denominators as stated,
\begin{equation}
S_6=\frac{80(45c+193)}{3c^{3}(5c+22)^{2}},\qquad
S_7=-\frac{2880(15c+61)}{7c^{4}(5c+22)^{2}},
\end{equation}
and at leading order in $1/c$ the whole tower collapses to a
one-parameter family, $S_r=\tfrac2r(-3)^{r-4}(2/c)^{r-2}+O(c^{-(r-1)})$,
whose Stieltjes transform is a single logarithm: the leading Virasoro
tower is the expansion of $-\tfrac{c^{2}}{162}\log(1+6t/c)$, one
Kummer motive with parameter $-6/c$, iterated against itself. This is
the algebraic channel of Section~\ref{sec:arithmetic} in its most
concrete form.

\subsection{Growth and resurgence}

The square root in \eqref{eq:riccati} places the branch points of $Q$ in
control of the asymptotics: $S_r\sim A\,\rho^{\,r}r^{-5/2}\cos(r\theta+\varphi)$
with $\rho=\sqrt{9\alpha^{2}+2\Delta}/2|\kappa|$. For the Virasoro family
$\rho(c)^2=(180c+872)/[c^{2}(5c+22)]$; the tower converges for $c$ beyond
the largest root $c_{*}\approx6.125$ of $5c^{3}+22c^{2}-180c-872$, so
every unitary minimal model sits in the divergent regime while $c=26$
converges. The Borel transform of the tower is entire; the field-theoretic
series it shadows is Gevrey-1 with the Koszul-dual saddle as its Borel
singularity. The shadow is the convergent projection of a divergent
theory.

\section{The landscape in detail}\label{sec:landscape}

The classification of Theorem~\ref{thm:classification} is only as good as
the computations that force it. This section records them, family by
family, at chain level.

\subsection{Heisenberg: the atom}

For $\mathcal{H}_k$ every structure of the paper is visible in closed
form. The bar chain space in degree $n$ has dimension $p(n-2)$, the
partition number, and the differential acts only through the double pole.
The degree-two computation
$d\bigl([\alpha|\alpha]\otimes\eta_{12}\bigr)=k\cdot\mathbf{1}$ exhibits
the curvature as the level; the degree-three cancellation is the whole
mechanism of Theorem~\ref{thm:selection} in miniature: on
$\alpha^{\otimes4}\otimes\eta_{12}\wedge\eta_{23}\wedge\eta_{34}$ only the
three adjacent divisors contribute, non-adjacent pairs dying by
$\eta\wedge\eta=0$ or by degree overflow, and $d^{2}$ telescopes to
$k^{2}-k^{2}=0$. The collision residue is $r(z)=k/z$; its
path-ordered exponential is the scalar $R$-matrix $R(z)=e^{k\hbar/z}$;
and the Koszul dual is the curved symmetric coalgebra of
Example~\ref{ex:heisenberg}. The genus-one block is the Fredholm
determinant of Section~\ref{sec:blocks}; the genus-$g$ block is the
string measure of Section~\ref{sec:physics}. Nothing in the general
theory is active here, and the general theory suffices in full.

\subsection{Affine algebras: two channels at work}

For $\widehat{\mathfrak{sl}}_2$ at genus two the curvature computation of
Theorem~\ref{thm:curvature} splits cleanly. The double-pole channel
contributes because $\partial_w\omega_j\neq0$ from genus two on, giving
$3k/4$ after the trace $\kappa^{ab}\kappa_{ab}=3$ and the normalisation
by $2\dvee=4$; the simple-pole channel closes a loop through the Killing
contraction $\sum_{c,d}f^{ac}_{\;\;d}f^{bc}_{\;\;d}=2\dvee\kappa^{ab}$
and contributes $\dim\fg/2=3/2$; the total is
$\kappa=3(k+2)/4$, vanishing at the critical level $k=-2$ where the two
channels cancel exactly. Nilpotency mixes the channels:
$\{d^{(0)},d^{(2)}\}+(d^{(1)})^{2}=0$ with both summands individually
nonzero --- a cancellation invisible in the abelian case, and the
precise content of Proposition~\ref{prop:bracket-fails}.

The Serre relations enter as degree-three bar cocycles, and the Jacobi
identity kills the quartic shadow: on the test vector
$(e,f,e,f)$ the three channels contribute
$\tfrac{2}{k+2}+0-\tfrac{2}{k+2}=0$. This is class $\mathsf{L}$
termination, witnessed rather than asserted. The bar cohomology, of
dimension $2n+1$ in degree $n$, is Chevalley--Eilenberg cohomology of
the loop algebra corrected by the Orlik--Solomon factor: for the Witt
algebra the mode-algebra count gives $\dim H^{1}_{\mathrm{CE}}=2$ where
the chiral answer is $1$, and the discrepancy is exactly the collapse of
non-consecutive collisions: the chiral bar complex is its own
object, with its own count.

\subsection{Level one and the unique class transition}

At level one the generic formula
$\kappa=\dim\fg\,(k+\dvee)/2\dvee$ fails for simply-laced $\fg$: the
value it predicts for $\widehat{\mathfrak{sl}}_2$ is $9/4$, the true
value is $1$. The mechanism is the Frenkel--Kac construction: the root
currents are composites of $\operatorname{rk}\fg$ Heisenberg fields, the
simple-pole channel dies, the cubic shadow vanishes, and the algebra
drops from class $\mathsf{L}$ to class $\mathsf{G}$ with
$\kappa=\operatorname{rk}\fg$. This is the unique crossing between
classes on the standard landscape, and it shows $\kappa$ is an invariant
of the algebra, not of a presentation: it must be computed, not read off
a formula.

\subsection{Symmetric orbifolds}

The modular characteristic is linear under symmetrisation:
$\kappa(\operatorname{Sym}^{N}\cA)=N\kappa(\cA)$, by four independent
mechanisms --- the twisted sectors carry positive weight and decouple
from the genus-one trace; the logarithm of the
Dijkgraaf--Moore--Verlinde--Verlinde product isolates the linear term;
the Hecke operator $T_N$ multiplies the vacuum obstruction by $N$;
and the bar complex of the orbifold splits into $N$ copies at genus
one. Four proofs of one linearity is the redundancy the subject
affords: every scalar statement in this paper is overdetermined.

\subsection{The $\beta\gamma$ system: three invariants split}

The $\beta\gamma$ system of weight $\lambda$ has
$\kappa=6\lambda^{2}-6\lambda+1$, the Mumford exponent; its collision
residue vanishes, the surviving ordered datum being the contact tensor
$\Theta=\beta\otimes\gamma-\gamma\otimes\beta$ with
$\Theta^{2}=0$ and $R=\mathrm{id}+2\pi i\,\Theta$; its Koszul dual is
the $bc$ system, the residue pairing exchanging the Weyl and Clifford
relations by an explicit four-term cancellation. On the weight family,
$S_3=0$ and $S_4=-5/12$ with all higher shadows vanishing --- the
class-$\mathsf{C}$ witness --- while on the induced Virasoro line at
$c(\lambda)$ the tower is infinite. The triple
$(p_{\max},k_{\max},r_{\max})=(1,0,4)$ is maximally split: pole order,
absorbed pole order, and shadow depth are three independent invariants,
and $\beta\gamma$ separates them.

\subsection{Bar dimensions and dual series}

The Hilbert series of the bar cohomologies close into algebraic
functions whose branch data is itself an invariant. Writing $h_n=\dim
H^{n}$ on $\mathbb{P}^{1}$:
\begin{center}
\begin{tabular}{lll}
\emph{algebra} & \emph{generating function} & \emph{discriminant}\\
\hline
$\mathcal{H}_k$ & $1/(1-x)$-type (partitions shifted) & $1$\\
$\beta\gamma$ & $\sqrt{(1+x)/(1-3x)}$ & $(1-3x)(1+x)$\\
$\widehat{\mathfrak{sl}}_2$ & $x(3-x)/(1-x)^{2}$ & $(1-3x)(1+x)$\\
$\Vir_c$ & $4x/\bigl(1-x+\sqrt{1-2x-3x^{2}}\bigr)^{2}$ & $(1-3x)(1+x)$\\
$\widehat{\mathfrak{sl}}_3$ & $4x(2-13x-2x^{2})/[(1-8x)(1-3x-x^{2})]$ &
$(1-8x)(1-3x-x^{2})$
\end{tabular}
\end{center}
The rank-one families share the branch divisor $(1-3x)(1+x)$ with
dominant growth $3^{n}=\dim\mathfrak{sl}_2$; the rank-two family grows
as $8^{n}=\dim\mathfrak{sl}_3$ with the golden-ratio-type factor
$1-3x-x^{2}$. Growth is governed by the dimension of the underlying Lie
algebra; the subdominant factor is a finer invariant, stable under
Drinfeld--Sokolov reduction at the level of branch charts. The Virasoro
dimensions are first differences of Motzkin numbers; the entropy
$h=\log(\text{growth rate})$ of the completed dual ascends strictly
along the $\cW$-tower, from $h(\cW_2)\approx0.567$ to the MacMahon value
$h(\cW_\infty)\approx0.872$: each additional generator weight adds
strictly positive completion entropy.

\subsection{The boundary of the landscape}

The Ising algebra is class $\mathsf{M}$ with an asymptotic,
Borel-summable tower ($S_4=40/49$, growth $\rho\approx12.5$): a
rational, solvable conformal field theory whose shadow expansion
diverges --- divergence of the tower is a statement about the
expansion, not about the theory. The Moonshine module has
$\kappa=c/2=12$, the double-pole channel carrying everything at
weight one, and
$\Delta=20/71\neq0$ places it in class $\mathsf{M}$; the Leech lattice
algebra at the same central charge is class $\mathsf{G}$ with
$\kappa=24$. The logarithmic triplet algebras $\cW(p)$ bound the landscape on the
non-semisimple side: $C_2$-cofinite with $2p$ simple modules and
provably outside finite free strong generation, they carry unbounded
Massey structure where the standard classes carry towers, and the
classification of this paper stops at their doorstep by
construction --- the class axioms presume the semisimple vacuum
geometry the triplets are built to violate. The critical-level
algebra $V_{-\dvee}(\fg)$ is uncurved,
non-Koszul, with trivialised monodromy and tripled degree-two
cohomology: every boundary phenomenon of the theory, in one example.

\subsection{Kinematics and dynamics}

The bar theory of an algebra splits into a kinematic and a dynamic
layer. The vacuum character alone fixes the kinematics: the growth of
the chain spaces, the entropy of the completed dual, the radius of the
Koszul series through the functional equation
\eqref{eq:character}. The operator product fixes the dynamics: the
differential, the tower, the class. Two algebras with the same
character and different products --- the Leech and Moonshine theories
at $c=24$ --- share every kinematic invariant and separate at the
first dynamic one, which is the cleanest justification for the
central place the tower occupies in this paper.

\subsection{The census}

The standard landscape, in the counting of this paper, comprises
twenty-one families: five Gaussian, two Lie-type, four contact, and
ten gravitational, with the critical stratum transversal to all. The
three pole-invariants, the shadow class, the modular characteristic,
and the conductor are tabulated in Appendix~\ref{app:census}; every
entry is the output of the two-channel computation of
Theorem~\ref{thm:curvature} and the recursion of
Theorem~\ref{thm:riccati}, and the table is closed under Koszul
duality, level reflection, and Drinfeld--Sokolov transport.

\section{Drinfeld--Sokolov and the $\cW$ quantum group}
\label{sec:DS}

\subsection{The escalator}

Quantised Drinfeld--Sokolov reduction along a principal nilpotent maps
class $\mathsf{L}$ to class $\mathsf{M}$: it manufactures gravity from
gauge theory. At chain level the mechanism is homotopy transfer: the
reduction contracts the bar complex along the BRST differential, and
the transferred $A_\infty$ operations are sums over planted binary
trees --- Catalan-many at each arity --- whose internal edges carry the
Drinfeld--Sokolov homotopy. The transfer sends the simple-pole
$r$-matrix of the affine algebra to the third-order kernel
$r(z)=(c/2)z^{-3}+2Tz^{-1}$ of the Virasoro algebra; it escalates the
pole by two, and with it the depth from $3$ to $\infty$. The coproduct
is transferred to its primitive part: a ghost-number argument
annihilates every group-like correction, which is the coalgebraic statement behind the
statement of Section~\ref{sec:physics} that gravity braids, and
braiding is all it does.

For $\cW_3$ the transferred data are explicit: the generator $W$ of
weight three has
$W_{(5)}W=c/3$ and
$W_{(1)}W=\tfrac{3}{10}\partial^{2}T+\tfrac{16}{22+5c}\Lambda$,
with $\Lambda$ the quasi-primary of \eqref{eq:S4}; the two-channel
decomposition $\kappa(\cW_3)=\tfrac{c}{2}+\tfrac{c}{3}=\tfrac{5c}{6}$
weights the generators by their Eulerian factors; the central-charge
conductor is $c+c^{!}=100$ with self-dual point $c=50$. The
non-principal reductions test the limits: for the Bershadsky--Polyakov
algebra $\cW^{k}(\mathfrak{sl}_3,f_{\mathrm{min}})$, with
$c(k)=-(2k+3)(3k+1)/(k+3)$, the reflection $k\mapsto-k-6$ gives the
exact conductor $c(k)+c(-k-6)=50$; the algebra is the first with mixed
shadow classes --- class $\mathsf{G}$ on its current line, class
$\mathsf{M}$ on its stress line --- and its modular characteristic is
the genus-one curvature of a non-principal reduction, the next
computation of the theory. For the $(3,2)$ nilpotent in
$\mathfrak{sl}_5$ the BRST differential contains the cubic ghost term
$c_{12}c_{23}b_{13}$: the dual object is an inhomogeneous Koszul
complex, and it governs the transport of duality along non-hook
reductions.

For hook-type nilpotents the duality transports along the reduction:
granting the compatibility of the reduction with the bar construction
edge by edge, the dual of $\cW^{k}(\mathfrak{sl}_N,f_{\lambda})$ is
$\cW^{-k-2N}(\mathfrak{sl}_N,f_{\lambda^{t}})$, the transpose
partition at the reflected level --- Koszul duality exchanging a
nilpotent orbit with its transpose, which is the affine shadow of the
duality of Barbasch and Vogan \cite{BV}.

\subsection{The quantum group of $\cW_{1+\infty}$}

The ordered theory of Section~\ref{sec:ordered} assigns to the
$\cW_{1+\infty}$ family a genuine quantum-group structure, and its
failure to be Hopf is a theorem. With parameter $\Psi$ and the scalar
Maulik--Okounkov $R$-matrix
$g(z)=\prod_i(z-h_i)/(z+h_i)$ \cite{MO}, the Drinfeld coproduct is
spectral,
\begin{equation}
\Delta_z\bigl(T(u)\bigr)=T(u)\otimes T(u-z),
\qquad
\Delta_z(T)=T\otimes1+1\otimes T
+\frac{\Psi-1}{\Psi}\,J\otimes J+z\,(1\otimes J)+\cdots,
\end{equation}
and the shift $u\mapsto u-z$ is forced by the conformal anomaly: the
obstruction class of Theorem~\ref{thm:anomaly} is exactly what a
$z$-independent coproduct would have to kill. The Miura cross-term
coefficient $(\Psi-1)/\Psi$ is independent of the spin --- verified
through spin six --- and reappears, with the same value, as the
obstruction cocycle
$[Q_{\alpha},\Delta_z^{\mathfrak{h}}]
=\tfrac{\Psi-1}{\Psi}\,e^{\alpha\cdot\varphi}\otimes
e^{\alpha\cdot\varphi}(1+O(z))$
to descending the free-field coproduct through the screening: the same
number measures the failure from both sides. An antipode exists at the
level of an involution, $S(J)=-J$,
$S(T)=-T+\tfrac{\Psi-1}{\Psi}:\!JJ\!:$, $S^{2}=\mathrm{id}$, but it is
a morphism of vertex algebras only at $\Psi\in\{1,2\}$, and the Hopf
axiom fails at every $\Psi$ with residual $-(\Psi-1)/z^{2}+zJ$: the
object is a chiral bialgebra and provably not a Hopf algebra. For
$\mathfrak{gl}_N$ the quantum determinant is central only for the
decreasing column ordering with spectral shifts $u-(i-1)\Psi$; the
ordering is rigid, the shift pattern forced.

\subsection{Genus one for the ordered theory}

At genus one the ordered bar complex of the affine algebra carries the
elliptic Knizhnik--Zamolodchikov--Bernard connection; its $B$-cycle
monodromy is $q^{2H}$ with $q=e^{\pi i/(k+\dvee)}$, and the elliptic
collision kernel decomposes through the Weierstrass functions,
$r^{(1)}=c_0\,\zeta_\tau+\sum_{n\geq1}\tfrac{(-1)^{n-1}}{n!}c_n\,
\wp^{(n-1)}$, with $B$-cycle defect $-2\pi i\,c_0$. The dichotomy of
the landscape reappears analytically: the curvature and the braiding
are entangled exactly when $c_0\neq0$, that is for the affine,
Virasoro and $\cW$ families, and decoupled for the free families. The
Belavin elliptic $r$-matrices, the Felder dynamical equation, and ---
at genus two --- a doubly-dynamical structure over the Siegel space
with the prime-form kernel, all arise as specialisations; at genus two
the correct object is a Hall--Drinfeld double rather than a Yangian,
the Manin-triple signature $(2,1)$ admitting no Lagrangian splitting.

\subsection{Seven faces of the $r$-matrix}

The binary collision residue \eqref{eq:rmatrix} admits seven
presentations in the literature of integrable systems: as the binary
component of the universal twisting morphism; as a line-operator
$r$-matrix; as the Laplace transform of a classical
$\lambda$-bracket; as the symbol of commuting sphere Hamiltonians; as
Drinfeld's classical $r$-matrix; as the kernel of the Sklyanin
bracket; and as the kernel of the Gaudin Hamiltonians, which arise
from it by taking residues,
$H_i=\sum_{j\neq i}\Omega_{ij}/(z_i-z_j)$, with the
Knizhnik--Zamolodchikov coupling $1/(k+\dvee)$. At genus zero the
seven faces coincide on the binary slice for a reason of
Grothendieck--Teichm\"uller type: the group $\mathrm{GRT}_1$ has no
generator in weight two, so no gauge separates them. At genus one the
quasi-period datum activates the higher weights and the seven elliptic
faces separate; the set of presentations becomes a torsor over
$\mathrm{GRT}_1(\bQ)$, which is the same torsor that indexes
splittings of the averaging map in Section~\ref{sec:ordered}. That the
ambiguity of quantisation and the ambiguity of presentation are one
group is forced; it is the associator doing its one job in two
costumes.

\section{The moduli of Koszulness and the fingerprint}
\label{sec:moduli}

\subsection{Koszulness as a moduli problem}

Koszulness is a space. For fixed $\cA$, consider
the functor of points whose values are triples: an associator
$\Phi$; a Maurer--Cartan deformation $\xi$ of the tower; and a
trivialisation of the cone of the quadratic comparison
$q_{\cA}$. At each finite Postnikov stage with framing this is an
open subfunctor of an affine Maurer--Cartan scheme, cut out by
Fitting ideals; the associator coordinate makes the whole object a
torsor over the Grothendieck--Teichm\"uller group, and Koszulness at
one associator implies it at every other. The moduli perspective
earns its keep twice. First, it organises the fourteen known
recognition tests --- among them the PBW collapse, the
$A_\infty$-formality, the monadicity criterion, the diagonal
concentration of factorization homology, the mixed-Hodge purity
test at the de Rham--Betti associator, and the Lagrangian
transversality test --- as charts of one atlas, each valid at its
own associator base point: ten of the tests are unconditional
biconditionals, three are one-directional, and the remainder carry
hypotheses. Second, it locates the failures of
Theorem~\ref{thm:non-koszul} as boundary strata: the Virasoro chart
closes non-circularly through the Feigin--Fuchs resolution, the Kac
determinant, and Hodge purity, while the admissible quotients sit
outside every chart, their Cartan obstructions visible from all of
them.

\subsection{The fingerprint}

The classification data assemble into the fingerprint
$\varphi(\cA)=(p_{\max},\,r_{\max},\,\chi,\,n_{\mathrm{strong}},\,
\text{coset})$: maximal generator pole, shadow depth, character,
number of strong generators, coset datum. On the standard landscape
the fingerprint together with the Koszul dual determines the bar
coalgebra up to isomorphism. The semion and its conjugate share
every scalar and differ by a twist: the twist is the datum beyond
the fingerprint, and the reconstruction clause carries a
residual-freeness hypothesis for exactly this reason. The coarse
projection of the fingerprint is the class of
Theorem~\ref{thm:classification}, and each projection has a
computable fibre: Virasoro and $\cW_3$ share a class and carry
conductors $13$ and $250/3$. What survives every projection is
$\kappa$; what survives none of them is the braiding.

\section{Complementarity and the conductor}\label{sec:conductor}

\subsection{The splitting}

Let $\cZ(\cA)$ be the centre local system on $\Mbar_{g}$ attached to a
chirally Koszul $\cA$ with Koszul dual $\cA^{!}$, and let
$\mathbf{C}_g(\cA)=R\Gamma\bigl(\Mbar_g,\cZ(\cA)\bigr)$.

\begin{theorem}[Complementarity]\label{thm:complementarity}
The Verdier involution $\sigma$ acts on $\mathbf{C}_g(\cA)$, and its
eigenspace decomposition splits the ambient cohomology into the
contributions of the two dual algebras,
\begin{equation}
\mathbf{C}_g(\cA)\;\simeq\;\mathbf{Q}_g(\cA)\,\oplus\,\mathbf{Q}_g(\cA^{!}):
\end{equation}
what one algebra counts as deformations, its dual counts as obstructions.
If in addition the pairing induced by Verdier duality is perfect and
anti-invariant, the two summands are complementary Lagrangians for a
symplectic structure of degree $-(3g-3+n)$, matching the dimension of
moduli. The scalar trace of the splitting is the \emph{conductor}
\begin{equation}\label{eq:conductor}
\varkappa(\cA)\;=\;\kappa(\cA)+\kappa(\cA^{!}),
\end{equation}
and its unconditional values on the standard landscape are:
\begin{equation}\label{eq:conductor-values}
\varkappa=0 \ \ (\text{Heisenberg, free fields, lattices, affine
under } k\mapsto-k-2\dvee),
\qquad
\varkappa(\Vir)=13 .
\end{equation}
\end{theorem}

\begin{proof}[Proof of the splitting]
The involution exists on the represented strict-flat locus: Verdier
duality on the Ran space composed with the centre comparison defines
$\sigma$ with $\sigma^{2}=\mathrm{id}$, and in characteristic zero the
idempotents $\tfrac12(1\pm\sigma)$ split any complex on which an
involution acts. The identification of the eigenspaces is the content:
the bar object of $\cA$ is the $j_{*}$-extension across the collision
divisors and lands in the $+1$ eigenspace, while the bar object of
$\cA^{!}$ is the $j_{!}$-extension, and the Kodaira--Spencer
anti-commutation $\mathbb{D}\circ\nabla_{\kappa}=-\nabla_{\kappa}
\circ\mathbb{D}$ --- contravariance of duality reversing infinitesimal
actions --- places it in the $-1$ eigenspace. Isotropy of each summand
and perfectness of the cross-pairing follow from anti-invariance of
the Verdier pairing when that pairing is perfect, which on the
standard landscape is verified through the PBW filtration: the
associated graded sees only the genus-zero differential, Koszulness
concentrates it in one row, and cohomology-and-base-change transports
perfectness across moduli. The scalar values are
Theorem~\ref{thm:curvature} plus arithmetic: $\kappa(V_k)+
\kappa(V_{-k-2\dvee})=0$ term by term, and
$\tfrac{c}{2}+\tfrac{26-c}{2}=13$.
\end{proof}

\begin{theorem}[Continuation of Theorem~\ref{thm:complementarity}]
\label{thm:complementarity-values}
The Virasoro value restates the arithmetic identity $c+c^{!}=26$, with
self-dual point $c=13$. Under the hypothesis that Drinfeld--Sokolov
reduction commutes with the bar construction, the principal
$\cW_N$ family carries
$\varkappa=(H_N-1)\bigl(4N^{3}-2N-2\bigr)$, in particular
$\varkappa(\cW_3)=250/3$ with self-dual point $c=50$; the central-charge
conductors $c+c^{!}=2\dim\fg$, $26$, $100$, and $4N^{3}-2N-2$ are exact
identities of rational functions of the level.
\end{theorem}

The two scalar conductors are related by the anomaly ratio: with
$K=c+c^{!}$ and $\varrho=\kappa/c=\sum_i 1/(m_i+1)$ over the exponents,
$\varkappa=\varrho\,K$. The ghost form of $K$ explains its values:
$K=-c_{\mathrm{ghost}}$ of the canonical BRST complex, and
$\sum_{j=2}^{N}2(6j^{2}-6j+1)=4N^{3}-2N-2$ recovers $26$ at $N=2$. That
the critical dimension of the bosonic string appears here as the Virasoro
conductor --- twice the self-dual point --- is the same cancellation,
met from the algebraic side.

\begin{conjecture}\label{conj:reflection}
$\Vir_c^{!}\simeq\Vir_{26-c}$ as chiral algebras. The scalar shadow of
this statement, $\varkappa=13$, is Theorem~\ref{thm:complementarity};
the fixed point of the reflection is the self-dual algebra at
$c=13$.
\end{conjecture}

\begin{conjecture}\label{conj:mukai}
Let $\widetilde H(K3,\bZ)=U^{4}\oplus E_8(-1)^{2}$ be the Mukai lattice,
of signature $(4,20)$, so that $2c_{+}=8$. The chiral algebra attached to
a $K3$ surface by the two-stage specialisation functor carries conductor
$\varkappa=8$, and the quantisation parameter obeys
$\hbar^{2}\varkappa=-1$, that is $\hbar^{2}=-\tfrac18$. The lattice
identity $2c_+=8$ is a theorem; the conjecture is the recognition of
the chiral source.
\end{conjecture}

For the Bershadsky--Polyakov algebra the central-charge conductor
$c(k)+c(-k-6)=50$ is an identity of rational functions; its modular
characteristic is the genus-one curvature of the minimal reduction,
and through the ratio law it carries the predicted conductor
$\varkappa=25/3$.

\subsection{The quantisation law}

The conductor fixes the quantum parameter. On the lattice side the
identity is arithmetic: for the reflective Borcherds data attached to
the Monster, $K3$ and Fake Monster lattices the pairs
$(\varkappa,\hbar^{2})$ are $(2,-\tfrac12)$, $(8,-\tfrac18)$ and
$(50,-\tfrac1{50})$, one law,
\begin{equation}
\hbar^{2}\,\varkappa\;=\;-1 :
\end{equation}
the square of the quantisation parameter is minus the inverse
conductor. On the chiral side the law is
Conjecture~\ref{conj:mukai} at the $K3$ row and a theorem of lattice
arithmetic everywhere it has been evaluated; it is the cleanest
statement of the principle that the deformation direction and the
duality pairing are inverse structures.

\subsection{Self-dual points}

The fixed locus of the Koszul involution is a subvariety of the space
of theories, and it is where the theory is most symmetric. The
self-dual points are the half-conductors: $c=13$ for the Virasoro
family, $c=50$ for $\cW_3$, $c=2N^{3}-N-1$ for $\cW_N$; among these
the Virasoro point is distinguished as the one attained at real
central charge --- the $\cW_3$ midpoint is reached only at
$k=-3\pm i$, so the Virasoro family alone carries its self-dual
theory on the physical line. On the fixed locus the complementarity
splitting balances, $\dim\mathbf{Q}_g(\cA)=\dim\mathbf{Q}_g(\cA^{!})
=\tfrac12\dim H^{*}$, and the conductor identity becomes the
equivariant signature of the involution. The three involutions of the
theory --- Verdier, the level reflection, and the braiding inversion
$q\mapsto q^{-1}$ --- coincide on this locus and only here; at
$c=13$ all three fix the same algebra.

\section{The genus expansion}\label{sec:genus-expansion}

\subsection{The scalar tower}

The genus expansion of $\cA$ is the sequence of traces of
$\Theta_{\cA}$ against the fundamental classes of $\Mbar_g$. Its first
term is unconditional; the higher terms factor through the Hodge bundle.

\begin{theorem}\label{thm:genus-expansion}
The genus-one trace is
$\operatorname{tr}_1\Ob^{\mathrm{def}}_{1,1}(\cA)=\kappa(\cA)\,\lambda_1$
in $H^{2}(\Mbar_{1,1})$; equivalently $F_1=\kappa/24$. Under the
$K$-theoretic comparison of the deformation class with
$\kappa\,\lambda_{-1}(\bE_g)$, the Hodge component at genus $g$ is
$(-1)^{g}\kappa\,\lambda_g$, and for algebras of uniform generator
weight the scalar free energies are
\begin{equation}\label{eq:Ahat}
F^{\mathrm{sc}}_g=\kappa\,\lambda_g^{\mathrm{FP}},
\qquad
\lambda_g^{\mathrm{FP}}
=\frac{2^{2g-1}-1}{2^{2g-1}}\,\frac{|B_{2g}|}{(2g)!},
\qquad
\sum_{g\geq1}F^{\mathrm{sc}}_g\,x^{2g}
=\kappa\Bigl(\frac{x/2}{\sin(x/2)}-1\Bigr),
\end{equation}
by Mumford's relation and the Faber--Pandharipande evaluation
\cite{Mum,FP}: the generating function is the Wick-rotated
$\widehat{A}$-genus, with radius of convergence $2\pi$ and universal
ratio $F_2/F_1=7/240$. The first six values of
$\lambda_g^{\mathrm{FP}}$ are
\[
\tfrac{1}{24},\quad
\tfrac{7}{5760},\quad
\tfrac{31}{967680},\quad
\tfrac{127}{154828800},\quad
\tfrac{73}{3503554560},\quad
\tfrac{1414477}{2678117105664000}.
\] For the $E_8$ lattice algebra
$F_2=7/720$, verified independently by the bar complex, by the
Siegel--Weil identity $\Theta^{(2)}_{E_8}=E_4^{(2)}$, and by direct
enumeration of the $30240$ pairs of roots.
\end{theorem}

The proof of Theorem~\ref{thm:genus-expansion} runs along the chain
Quillen anomaly, Grothendieck--Riemann--Roch, Faber--Pandharipande:
the curvature of the determinant line is $-2\pi i\,c_1(\bE)$; the
Chern character of the Hodge bundle is
$g+\sum_{j\geq1}\frac{B_{2j}}{(2j)!}\,\kappa_{2j-1}$ on the
universal curve; the $\psi$-integrals
$\int_{\Mbar_{g,1}}\psi^{2g-2}\lambda_g$ evaluate the numbers; and
Mumford's identity $\lambda_g^{2}=0$ resums the even Todd series to
the sine. Independent confirmations come from the family index
theorem and, for the free families, from the operator theory of
Section~\ref{sec:blocks}.


\subsection{The multi-weight correction}

\begin{theorem}\label{thm:multiweight}
The scalar formula \eqref{eq:Ahat} is exact only on the uniform-weight
lane, and the failure beyond is explicit. For $\cW_3$ at genus two, the
sum over the seven stable graphs of $\Mbar_2$, with channel assignments
weighted by the two-field propagators $\eta^{TT}=2/c$, $\eta^{WW}=3/c$
and the $\bZ/2$ selection rule killing odd-$W$ vertices, gives
\begin{equation}\label{eq:cross}
F_2(\cW_3)
=\frac{7\kappa}{5760}
+\delta F_2^{\mathrm{cross}},
\qquad
\delta F_2^{\mathrm{cross}}
=\frac{3}{c}+\frac{9}{2c}+\frac{1}{16}+\frac{21}{4c}
=\frac{c+204}{16\,c},
\end{equation}
the four terms being the sunset, theta, bridge-loop and barbell
graphs --- the remaining three of the seven stable graphs of
$\Mbar_2$ carry an odd number of $W$-legs at some vertex and vanish
by the $\bZ/2$ symmetry --- and the bridge-loop contribution $1/16$
is a topological residue independent of $c$. This enumeration is the
proof: seven graphs, three killed by parity, four evaluated by the
propagators $\eta^{TT}=2/c$, $\eta^{WW}=3/c$ against the structure
constants of $\cW_3$: the sum is
$\tfrac3c+\tfrac9{2c}+\tfrac{21}{4c}=\tfrac{51}{4c}$, and
$\tfrac{51}{4c}+\tfrac1{16}=\tfrac{c+204}{16c}$. The corrections grow with genus faster than the scalar part
--- the ratios of cross to scalar at $g=2,3,4$ are
$O(1)$, $42/31$, $9184/381$ --- so the $\widehat{A}$-tower is the
leading term, not the sum, of the free energy of an interacting
multi-weight algebra. Free-field algebras are the exact multi-weight
exception: all their cross-channels vanish identically.
\end{theorem}

The corrections themselves obey a universal law: for the principal
$\cW_N$ family the gravitational part of the genus-two correction is
\begin{equation}
\delta F_2^{\mathrm{grav}}(\cW_N)
=\frac{(N-2)(N+3)}{96}
+\frac{(N-2)(3N^{3}+14N^{2}+22N+33)}{24\,c},
\end{equation}
vanishing at $N=2$ as it must, exact at $N=3$, and growing as
$N^{2}/96$ in the planar limit --- the 't~Hooft scaling emerging from
the stable-graph combinatorics alone.

The two theorems together delimit what a scalar can know. The genus
tower of any chiral algebra begins with a universal multiple of the
$\widehat{A}$-genus; the deviation from universality is a graph
calculus with its own arithmetic --- the $\cW_4$ corrections leave the
field $\bQ(c)$ --- and it, not the scalar, dominates asymptotically.

\subsection{Resurgence of the scalar tower}

The generating function \eqref{eq:Ahat} is meromorphic with simple
poles at $x=2\pi n$, and its resurgent structure is completely
explicit: the large-genus asymptotics
$F_g\sim2\kappa\,(2\pi)^{-2g}$ are saturated by the first pole; the
instanton action is $(2\pi)^{2}$, universal across the landscape; the
Stokes constants are $(-1)^{n}\,4\pi^{2}n\,\kappa\,i$; and the Borel
transform is entire, so the tower is Borel-summable with optimal
truncation at $g\approx\pi/\rho$. Set against the $(2g)!$ growth of
the string perturbation series this is the cleanest formulation of
the paper's analytic moral: the modular trace projects a factorially
divergent theory onto a convergent shadow, and the projection loses
exactly the non-perturbative sectors, whose location it nevertheless
records in its Stokes data. The Gaussian matching completes the
picture: on the scalar lane the free energies coincide with those of
the Gaussian matrix model at $N^{2}=\kappa$, by Harer--Zagier \cite{HZ}; the
shadow of any chiral algebra is a Gaussian theory with coupling
$\kappa$, and everything beyond the shadow --- the graph calculus
of Theorem~\ref{thm:multiweight} --- is precisely what distinguishes
an interacting algebra from its Gaussian projection.

\section{Gaudin, Bethe, and the spectral chain}\label{sec:gaudin}

The integrable structure of a chiral algebra is the ordered theory read
spectrally, and it descends from $\Theta_{\cA}$ along one chain:
\begin{equation}
\Theta_{\cA}\;\longrightarrow\;r(z)\;\longrightarrow\;R(z)
\;\longrightarrow\;T(u)\;\longrightarrow\;Q(u):
\end{equation}
the Maurer--Cartan element projects to the classical $r$-matrix, whose
path-ordered exponential is the quantum $R$-matrix, whose trace is the
transfer matrix, whose Baxter operator carries the spectrum.

The first arrow is Theorem~\ref{thm:mc} at degree two. The second is
the Kohno monodromy of the shadow connection: flatness of
$\nabla=d-\sum r_{ij}(z_i-z_j)\,dz_{ij}$ is the classical Yang--Baxter
equation, which is the Arnold relation \eqref{eq:arnold} contracted
with the Casimir; the quantum Yang--Baxter equation for the monodromy
follows by path-ordering. The third arrow is the RTT trace, and
commutativity of transfer matrices is the geometric statement that
disjoint collision cycles commute. The fourth is representation
theory on the Baxter locus: the singular vector
$w_\lambda=\lambda(v_-\otimes v_\lambda)-v_+\otimes fv_\lambda$ is a
highest-weight vector exactly at $b=a-(\lambda+1)/2$, and the
$TQ$-relation $\Lambda Q=aQ(u-i)+dQ(u+i)$ follows from the resulting
short exact sequence of prefundamentals.

Bethe's equations close the circle: the stationarity of the shadow
Yang--Yang potential reproduces
$\prod_{j\neq i}\frac{u_i-u_j+i}{u_i-u_j-i}
=\bigl(\frac{u_i+i/2}{u_i-i/2}\bigr)^{L}$, with energies
$E=\tfrac{L}{4}-\tfrac12\sum_a(u_a^{2}+\tfrac14)^{-1}$, and the Bethe
vectors $B(u_1)\cdots B(u_M)\lvert\Omega\rangle$ live in the ordered
tensor product: the ordered bar complex is the chain-level refinement
of the algebraic Bethe ansatz, the symmetric shadow retaining the
spectrum $(\Lambda,Q)$ and the ordered kernel the wavefunctions. The
Gaudin Hamiltonians $H_i=\sum_{j\neq i}\Omega_{ij}/(z_i-z_j)$ are the
residues of the second arrow, with the Knizhnik--Zamolodchikov coupling
$1/(k+\dvee)$; their diagonalisation is the semiclassical face of the
same chain.

\section{Conformal blocks, Verlinde numbers, and the analytic layer}
\label{sec:blocks}

\subsection{Blocks from the pointed bar complex}

The zeroth cohomology of the bar complex with module insertions
computes derived coinvariants; the comparison with the sheaf of
conformal blocks of Tsuchiya--Ueno--Yamada \cite{TUY} holds given the integrable
truncation, the flat connection, and the sewing and determinant-line
data. On the integrable
quotient of $\widehat{\mathfrak{sl}}_2$ at level $k$ the identification
is complete enough to recover the Verlinde formula \cite{Verlinde} from chiral
homology: with $n=k+2$,
\begin{equation}
Z_g(k)=\sum_{j=1}^{n-1}S_{0j}^{\,2-2g},
\qquad
Z_g=\frac{n^{g-1}}{2^{g-1}}\sum_{j=1}^{n-1}
\csc^{2g-2}\!\Bigl(\frac{\pi j}{n}\Bigr),
\end{equation}
a polynomial in $n$ of degree $3g-3$; the closed forms begin
$Z_2=n(n^{2}-1)/6$ and $Z_3=n^{2}(n^{2}-1)(n^{2}+11)/180$, and the
leading coefficient is $\zeta(2g-2)/(2^{g-2}\pi^{2g-2})$, which is the
source of the Bernoulli numerators of Theorem~\ref{thm:irregular}. The
handle operator is $S_{0j}^{-2}$; the genus-one count $Z_1=k+1$ is
obtained three independent ways --- the $S$-matrix, the Zhu algebra \cite{Zhu}
$U(\mathfrak{sl}_2)/(e^{k+1},f^{k+1})$, and the degree-zero ordered
centre. At level one the interacting sewing kernel is diagonalisable in
closed form, with eigenvalues $\{q,q^{2},q^{2}\}$ on the three-state
block --- the first fusion computation of the analytic layer beyond the
free families.

\subsection{Characters from the bar resolution}

The classical character formulas are bar cohomology. Filtering the
bar complex of $\widehat{\fg}_k$ by the Weyl group, the first page is
a sum of Verma characters over affine Weyl chambers, the differential
is composed of the singular-vector embeddings located by the Kac
determinant, and the alternating cancellation is the
Bernstein--Gelfand--Gelfand resolution run inside the bar complex:
the Weyl--Kac formula
\begin{equation}
\operatorname{ch}L(\lambda)
=\frac{\sum_{w}\varepsilon(w)\,e^{w(\lambda+\rho)-\rho}}
{\prod_{\alpha>0}\bigl(1-e^{-\alpha}\bigr)^{\operatorname{mult}\alpha}}
\end{equation}
drops out of the collapse, and the denominator identities of the free
theories are the statement that their bar complexes are resolutions.
At admissible fractional level the same mechanism runs with screening
operators as the differentials: the Wakimoto free-field realisation \cite{Wakimoto}
resolves the vacuum module away from the Kac--Kazhdan hyperplanes,
its BRST cohomology concentrating in one degree \cite{Arakawa}, and the
Kac--Wakimoto characters collapse, by the Kac--Wakimoto theorem \cite{KW}, to a
two-term Weyl sum. Free-field realisations are, in this paper's terms, chains
of quasi-isomorphisms between bar complexes; the screening charges
are the connecting differentials; and the classical character theory
of affine algebras is the shadow of Theorem~\ref{thm:pbw} on
characters.

\subsection{The analytic layer}

The algebraic sewing of Theorem~\ref{thm:mc} acquires an analytic
completion under growth hypotheses, and for the free families the
completion converges.

\begin{theorem}\label{thm:sewing}
Suppose the graded dimensions of $\cA$ grow subexponentially,
$\dim\cA_n\leq Ce^{\alpha\sqrt n}$, and the structure constants
polynomially. Then the $q$-damped pair-of-pants operator is
Hilbert--Schmidt for every $0<q<1$, and every closed stable amplitude
with at least two sewn channels is trace class. For the rank-$r$
Heisenberg algebra the one-particle sewing operator $T_q$ has
eigenvalues $q^{n}$, trace norm $q/(1-q)$, and Fredholm determinant
\begin{equation}
\det(1-T_q)\;=\;q^{-1/24}\,\eta(q),
\end{equation}
so the genus-one block is $\eta(\tau)^{-r}$ by second quantisation,
and the genus-$g$ block is the string measure of
Section~\ref{sec:physics}. For lattice algebras the block acquires the
factor $\Theta_\Lambda(\Omega)$, and at genus two the theta factor
separates isospectral lattices beyond the reach of genus one.
\end{theorem}

\begin{proof}
The growth hypotheses bound the matrix elements of the damped
pair-of-pants by a summable Gaussian in the weights, giving the
Hilbert--Schmidt property; a closed amplitude with two sewn channels
is a product of two Hilbert--Schmidt operators, hence trace class.
For the Heisenberg algebra, Wick's theorem reduces the second
quantisation to the one-particle Bergman space, where the sewing
operator is multiplication by $q^{n}$ and the determinant is the
Euler product of $\eta$.
\end{proof}

The Dedekind eta function is an output here: it is the Fredholm
determinant of the sewing operator. This is the sense in which the
analytic theory is a layer over the same graph sums, and not a second
theory.

\subsection{The critical level}

The boundary stratum $\mathsf{FF}$ deserves its own statement. At
$k=-\dvee$: the curvature vanishes ($\kappa=0$), so the bar complex is
uncurved at every genus; the monodromy trivialises,
$e^{2\pi i/(k+\dvee)}$ degenerating; Koszulness fails, the cohomology
spreading off the diagonal; and the degree-zero bar cohomology is the
Feigin--Frenkel centre, the algebra of functions on opers for the
Langlands dual group \cite{FF}. Under the comparison of the bar
complex with the de Rham complex of the space of opers on the formal
disc, the higher cohomology is $\Omega^{n}(\mathrm{Op}_{{}^{L}\fg}(D))$
--- a statement on the formal disc, at the threshold where the
localisation theory of Frenkel--Gaitsgory begins. The level reflection $k\mapsto-k-2\dvee$ is the unique
affine involution fixing the critical level and compatible with
$\kappa$; on quantum parameters it is $q\mapsto q^{-1}$; and its fixed
stratum is where the Langlands-dual geometry replaces the finite
closed sector of Theorem~\ref{thm:H}. At rank one the critical centre is periodic: the formal-disc
Hochschild cohomology of $V_{-2}(\mathfrak{sl}_2)$ is the free
graded-commutative algebra $\Lambda(P)\otimes\bC[\Theta]$ with
$\deg P=3$, $\deg\Theta=4$ --- four-periodic, the algebraic germ of
the modularity that the paramodular forms of
Section~\ref{sec:moonshine} exhibit globally. Everything singular
about the critical level is the statement that flatness and duality
degenerate together.

\section{The genus-two kernel}\label{sec:genus-two}

Genus two is where the modular theory first outruns every genus-one
structure, and each of its new phenomena is computable.

The propagator is built from the Szeg\H{o} kernel
$S_2(z,w)=d\log E(z,w)$ with quasi-periodicity
$S_2(z+B_\alpha,w)=S_2(z,w)-2\pi i\,\omega_\alpha(z)$: the monodromy
defect, scalar at genus one, is function-valued at genus two, and
through it the level enters the curvature directly
(Theorem~\ref{thm:curvature}). The flat connection acquires three
modular directions, one for each entry of the period matrix, with the
genus-two heat equation
$\partial_{\Omega_{\alpha\beta}}\theta=(2\pi i)^{-1}
\partial_{z_\alpha}\partial_{z_\beta}\theta$ coupling them; the
twisted local system on the configuration space of a genus-two curve
has $\dim H^{1}=12$ at generic monodromy, by the Euler-characteristic
count $\chi=r(2-2g-s)$. The quantum object attached to the
two-holed geometry is a Hall--Drinfeld double: the Manin triple has
signature $(2,1)$, and the two-dimensional positive part against the
one-dimensional negative part replaces the Lagrangian splitting of the
Yangian --- genus two quantises a double, and the doubling is forced
by the signature.

On the automorphic side the genus-two kernel reaches what genus one
provably organises differently: the Siegel theta series separates
$E_8\oplus E_8$ from $D_{16}^{+}$; the two-variable sewing kernel
carries an off-diagonal Schur complement whose spectral decomposition
reaches central values of twisted $L$-functions through the Igusa form
$\Phi_{10}=\Delta_5^{2}$; and the second-quantised $K3$ elliptic genus
$1/\Phi_{10}$ sits at the boundary of the classification, where the
Borcherds algebras of Section~\ref{sec:moonshine} begin. The
genus-two kernel is the smallest window through which the arithmetic
of Section~\ref{sec:arithmetic}, the automorphic boundary, and the
gravitational classes are all visible at once.

\section{The bulk}\label{sec:bulk}

\subsection{The derived centre as the closed sector}

The bar complex of Section~\ref{sec:bar} lives on configurations in a
complex curve crossed with an ordered real line: the holomorphic factor
carries the residues, the topological factor the deconcatenation. This is
the geometry of a three-dimensional holomorphic--topological theory on
$X\times\bR$, encoded operadically by the two-coloured Swiss-cheese operad
$\mathrm{SC}^{\mathrm{ch,top}}$, whose closed colour is $\FM_k(\bC)$,
whose mixed part is $\FM_k(\bC)\times\Conf^{<}_m(\bR)$, and whose
open-to-closed part is empty. The algebra over this operad is the pair assembled from the boundary
algebra and its centre; the bar coalgebra is the resolution by which
the centre is computed:
\begin{equation}
\cZ^{\mathrm{der}}_{\mathrm{ch}}(\cA)
=R\mathrm{Hom}_{\cA\otimes\cA^{\mathrm{op}}}(\cA,\cA)
=\ChirHoch^{\bullet}(\cA,\cA),
\end{equation}
the chiral Hochschild complex, carrying cup product and brace operations
extracted from insertion at boundary strata of $\FM_{k+1}$, hence a
homotopy Gerstenhaber ($E_2$) structure: the chiral form of Deligne's
conjecture.

\begin{theorem}[Terminality]\label{thm:terminality}
The pair
$\bigl(\cZ^{\mathrm{der}}_{\mathrm{ch}}(\cA),\,\cA\bigr)$, with the brace
action of the centre on the algebra, is the terminal local open--closed
pair over $\cA$: every $\mathrm{SC}^{\mathrm{ch,top}}$-pair
$(\cB,\cA)$ admits a unique morphism of pairs
$\cB\to\cZ^{\mathrm{der}}_{\mathrm{ch}}(\cA)$, given on elements by
$b\mapsto(-1)^{|b|}\mu_{1;n}(b;-)$. Every local closed sector acting on
a chiral algebra factors uniquely through its derived centre.
\end{theorem}

\begin{proof}
The codimension-one identities of the Swiss-cheese operad force the
assignment to be a chain map compatible with the braces; evaluation on
the boundary algebra forces its value. Existence and uniqueness are the
two halves of the same computation, and terminality follows.
\end{proof}

Theorem~\ref{thm:terminality} is holography as a universal property. The
third dimension is constructed: it is the normal collar that the
centre construction manufactures. The theorem is also exactly as strong
as such a statement can be:

\begin{proposition}[Separation]\label{prop:separation}
A decoupled closed sector with zero boundary coupling acts trivially
on every entry of the algebraic data
\[
\bigl(\cA,\ \cA^{i},\ \cA^{!},\
\cZ^{\mathrm{der}}_{\mathrm{ch}}(\cA),\ r(z),\ \Theta_{\cA}\bigr),
\]
while enlarging the bulk observables of any holomorphic--topological theory
that carries it. A physical bulk is therefore specified by the
algebraic data together with an open--closed comparison map
identifying $\cZ^{\mathrm{der}}_{\mathrm{ch}}(\cA)$ with its boundary
observables: the comparison map is an independent coordinate of the
theory.
\end{proposition}

\subsection{The annulus and the trace}

Evaluation on the circle produces the trace: by excision over the
two-interval cover, the factorization homology of the annulus is the
two-sided bar construction
$\cA\otimes^{L}_{\cA\otimes\cA^{\mathrm{op}}}\cA$ --- Hochschild
homology --- and the identification is Morita invariant. The cyclic
structure of the trace is the datum that thickens graphs to ribbon
graphs, and with it the genus expansion acquires its 't~Hooft form:
a matrix realisation of the algebra weights each stable graph by
$N^{\chi}$, and the planar limit of Section~\ref{sec:genus-expansion}
is literal. The closed-sector trace, the cyclic symmetry of
correlators, and the large-$N$ counting are one structure, met three
times.

\subsection{Topologisation}

The centre is holomorphic in the curve direction; it becomes topological
exactly when translation is inner.

\begin{theorem}\label{thm:E3}
Let $\cA=V_k(\fg)$, $k\neq-\dvee$, with Sugawara vector
$T=\tfrac{1}{2(k+\dvee)}\sum:\!J^{a}J_{a}\!:$. On the BRST cohomology of
the derived centre, $T=[Q,G]$ for an explicit $G$; translations are
therefore exact, the factorization algebra is locally constant on
cohomology, and Dunn additivity upgrades the $E_2$-structure to
$E_3$: the derived centre of an affine algebra at non-critical level is
a three-dimensional topological object. At the critical level the
Sugawara construction diverges and the topologisation collapses onto the
commutative Feigin--Frenkel centre. For class~$\mathsf{M}$ the
statement holds on cohomology and on weight-completed models; on the
original complex it is equivalent to an $A_\infty$-coherent lift of
$[m,G]=\partial_z$, the remaining step.
\end{theorem}

\begin{proof}
$\partial_z\mathcal{O}=[T,\mathcal{O}]=[Q,[G,\mathcal{O}]]$ makes
translation exact; a factorization algebra with exact translations is
locally constant on cohomology; Lurie's recognition principle converts
local constancy into a topological $E_1$-factor transverse to the
holomorphic $E_2$; Dunn additivity multiplies them.
\end{proof}

\subsection{Three cohomologies}

Three cohomology theories shadow the centre, on three domains: the
curve-level chiral Hochschild complex of this section; the Hochschild
cohomology of a chosen mode algebra on the formal disc; and the
Gelfand--Fuchs continuous cohomology of the positive modes. The
comparison maps run from the curve to the disc --- evaluate at a
colliding configuration and set the spectral variables to zero ---
and from the disc to Gelfand--Fuchs by antisymmetrisation against an
invariant functional; each map forgets the data proper to its source,
the spectral variables in the first case and the full associative
structure in the second. The three theories agree in low degrees on
the benchmarks and diverge exactly where they should: the
Gelfand--Fuchs side is unbounded where the chiral side is finite, and
the difference is measured by the cone of the comparison, which is
itself computable.

\subsection{Support of the centre}

\begin{theorem}\label{thm:H}
Suppose given a family datum for $\cA$: a filtered model of
$\ChirHoch^{\bullet}(\cA)$ together with a strong deformation retract
onto a complex supported in a finite set $S$ of degrees. Then the
support of $\ChirHoch^{\bullet}(\cA)$ lies in $S$. The benchmarks, due
to Bakalov--De Sole--Kac \cite{BDSK}, are: for $\Vir_c$, support
$\{0,2,3\}$ with one-dimensional groups --- the Poincar\'e polynomial is
$1+t^{2}+t^{3}$, and cup product with the anomaly class is an
isomorphism $H^{0}\to H^{2}$; for the rank-one even superboson, support
$\{0,1\}$ with dimensions $(2,1)$. When a perfect pairing of degree $d$ aligns
the two charts, the support is palindromic:
$\ChirHoch^{n}(\cA)\cong\ChirHoch^{d-n}(\cA^{!})^{\vee}\otimes\omega_X$
and $P_{\cA}(t)=t^{d}P_{\cA^{!}}(t^{-1})$ --- at $d=2$ the
deformations of $\cA$ are the obstructions of $\cA^{!}$, exchanged by
duality. For the affine, Virasoro and
$\cW$ families the retract from collision geometry is the hypothesis
under which the theorem operates. At the critical level the support is
infinite: the centre is the algebra of functions on opers, with
Hilbert series
$\prod_{i=1}^{r}\prod_{m\geq0}(1-q^{\,d_i+1+m})^{-1}$. Finiteness of
the closed sector is a property of non-degenerate vacua, and it ends
exactly where the Langlands-dual geometry begins.
\end{theorem}

\section{Arithmetic of the tower}\label{sec:arithmetic}

\subsection{Two channels}

The shadow tower meets arithmetic through two channels of opposite
transcendence, and the asymmetry between them is structural.

The \emph{algebraic channel} is Theorem~\ref{thm:riccati}: shadow
extraction is a depth-one residue against Arnold forms --- no iterated
integrals occur --- so the tower of any single-line family is rational.
For the Virasoro family more is true: at leading order in $1/c$ the
recursion closes on a single plethystic primitive, and the entire tower
is the iterated self-extension of one Kummer motive with parameter
$-6/c$; the leading Laurent coefficients are
$A_r=8(-6)^{r-4}/r$, with the subleading tier
$B_r=-A_r\bigl(\tfrac{22}{5}+\tfrac{(r-4)(r-5)}{18}\bigr)$ in closed
form --- the Lee--Yang constant $\tfrac{22}{5}$ governing the entire
subleading stratum. Infinite algebraic depth, zero transcendence. The
associator, with its multiple zeta values, lives in the kernel of the
averaging map and never reaches the tower.

The \emph{arithmetic channel} opens for lattice algebras, whose graph
amplitudes factor through theta series in a finite-dimensional Hecke
module. The tower then resolves the constrained Epstein zeta function
one Hecke constituent per degree: degree two reads the constant term,
degree three the Eisenstein part, degree $3+j$ the $j$-th cusp
eigenform, and termination is the dimension of the space of cusp forms.
For the Leech lattice the quartic obstruction \emph{is} the Ramanujan
$L$-function $L(s,\Delta_{12})$. Finite depth, periods of full
transcendence. The resolution is worked in rank: for $\bZ$ the tower
stops at once; for $E_8$ there is one Eisenstein and one cusp
constituent; for the Leech lattice the theta series is
$E_{12}-\tfrac{65520}{691}\Delta_{12}$, the Ramanujan congruence
$\tau(n)\equiv\sigma_{11}(n)\bmod 691$ reading off the constant, and
the tower terminates at depth four because
$\dim S_{12}(\mathrm{SL}_2(\bZ))=1$.

\subsection{Irregular primes}

\begin{theorem}\label{thm:irregular}
Let $Z_g(k)$ be the genus-$g$ Verlinde number of $\widehat{\mathfrak{sl}}_2$
at level $k$, a polynomial in $n=k+2$ of degree $3g-3$:
$Z_g=n^{g-1}\sum_{j=1}^{n-1}\csc^{2g-2}(\pi j/n)/2^{g-1}$. Its leading
coefficient is $(-1)^{g}2^{g-1}B_{2g-2}/(2g-2)!$, by the Hurwitz
expansion of the cosecant against Euler's evaluation of $\zeta(2g-2)$.
Consequently the Kummer-irregular primes enter the Verlinde polynomials
through Bernoulli numerators: $691$ divides the numerator of
$[n^{18}]Z_7$ and $3617$ divides that of $[n^{24}]Z_9$, each prime
first appearing at genus $g(p)$ with $p>2g-2$, so that the factorial
factor is a $p$-adic unit. In the computed range $r\leq11$ the cleared
numerators of the Virasoro tower contain no Kummer-irregular prime:
the two channels are disjoint where both have been computed.
\end{theorem}

\begin{proof}
The cosecant power sum is polynomial in $n$ by the Eisenstein--Cauchy
argument; its two boundary tails give
$2\zeta(2g-2)(n/\pi)^{2g-2}$ by Hurwitz's expansion, and Euler's
evaluation converts $\zeta(2g-2)$ into $B_{2g-2}$. Witten's
symplectic volume of the moduli of flat connections \cite{WittenVol} provides the
independent second derivation of the same leading coefficient.
\end{proof}

\subsection{The Galois layer}

The paramodular forms of the genus-two kernel carry Galois theory,
and the first floor is computed. The spinor $L$-function of the
Saito--Kurokawa lift attached to the Igusa form has Hecke eigenvalues
$\lambda_p=a_p(f_{18})+p^{8}+p^{9}$ with $f_{18}=E_6\Delta$, its
degree-four factor $p^{34}$; the associated deformation ring is
non-reduced,
$\bZ_\ell[[x_1,x_2,x_3]]/(x_1^{2},x_1x_2,x_1x_3)$, the nilpotent
direction cut out by a cup-square in Galois cohomology --- the
arithmetic counterpart of a first-order deformation that obstructs
itself, which is the Galois-theoretic face of the quartic
obstructions this paper computes on the chiral side. The matching of
the two obstruction theories across the Igusa kernel is the
arithmetic form of the modular Koszul programme.

\subsection{Dirichlet lifts of the sewing operator}

The sewing spectrum organises into Dirichlet series, and the free and
gravitational families lift differently:
\begin{equation}
\begin{aligned}
S_{\mathcal{H}}(u)&=\zeta(u)\,\zeta(u+1),\qquad
S_{\Vir}(u)=\zeta(u+1)\bigl(\zeta(u)-1\bigr),\\
S_{\cW_N}(u)&=\zeta(u+1)\Bigl((N{-}1)\zeta(u)-\textstyle\sum_{j<N}H_j(u)\Bigr),
\end{aligned}
\end{equation}
the Heisenberg lift carrying the full Euler product of
$\zeta(u)\zeta(u+1)$ --- the logarithm of the Fredholm determinant of
Section~\ref{sec:blocks} expanded in divisor sums --- and each
$\cW$-generator subtracting its own harmonic correction: the Miura
transformation acts on the arithmetic side as a divisor-sum defect.
For even unimodular lattices the depth of the Hecke resolution is
$3+\dim S_{r/2}$: depth two for $\bZ$, three for $E_8$, four for the
Leech lattice, whose cusp constituent is the Ramanujan form.

\subsection{The reach of genus one}

Three theorems fix the reach of the genus-one theory. Rankin--Selberg
unfolding of any genus-one integral against an Eisenstein series
retains exactly the constant term of the zeta function; the modular
Maurer--Cartan constraint lives in the pluriharmonic direction, which
Weil--Petersson geometry projects out; and the genus-one sewing
operator is Fock-diagonal, so every two-variable Dirichlet series it
produces collapses to one variable. Genus-one Maurer--Cartan data is
integrability: it fixes the couplings. The location of zeros is a
genus-two phenomenon: the sewing kernel of genus two carries an
off-diagonal Schur complement, and through the Igusa cusp form it
reaches the central values of twisted $L$-functions --- the route the
next paper takes.

\section{Deformation quantization and the $E_3$ identification}
\label{sec:E3}

\subsection{Quantization as Maurer--Cartan theory}

A coisson algebra --- a commutative chiral algebra with a Lie${}^{*}$
bracket --- is quantized by lifting a Maurer--Cartan element order by
order in $\hbar$ through its chiral deformation complex; the local
input is formality on the disc, the global obstructions live in the
modular directions, and the classification of quantizations is the
gauge groupoid of the complex. Level quantization is the first
example: for the affine algebra the obstruction lattice is
$H^{3}(\fg;\bZ)$, and integrality of the level is the integrality of
the Maurer--Cartan class.

The boundary form of quantization is a theorem of Koszul type. For a
potential $F$, the cofree coalgebra on the graph of $dF$ carries a
canonical coderivation, and the coderivation squares to zero exactly
on the derived critical locus: for $F=x^{2}$ the square evaluates to
$6c^{2}$ on the quartic cogenerator, and the terminal square-zero
quotient is the Koszul complex of the equations. On that quotient the
Hochschild cohomology is the algebra of functions on the shifted
cotangent of the derived zero locus,
$HH^{\bullet}(A_F)\simeq\mathcal{O}\bigl(T^{*}[-1]Z_F^{\mathrm{der}}
\bigr)=\mathcal{O}\bigl(d\mathrm{Crit}\bigr)$ ---
Landau--Ginzburg theory as the simplest instance of the principle
that curvature, unabsorbed, cuts out geometry.

\subsection{The identification with Chern--Simons theory}

For $\cA=V_k(\fg)$ at non-critical level the derived centre is small
and completely computable on the formal disc:
\begin{equation}
\cZ^{\mathrm{der}}_{\mathrm{ch}}\bigl(V_k(\mathfrak{sl}_2)\bigr)
\;=\;\bC\,\oplus\,\mathfrak{sl}_2[-1]\,\oplus\,\bC[-2],
\end{equation}
with vanishing cup product and Gerstenhaber bracket and with the
$P_3$-bracket
\begin{equation}
\{X,Y\}\;=\;\frac{(X,Y)}{k+\dvee}
\end{equation}
as its only structure: the entire local closed sector of the
Wess--Zumino--Witten model is the invariant form divided by the
shifted level. On the other side, perturbative Chern--Simons theory
for $\fg$ produces a filtered $E_3$-algebra of observables whose
deformation space is $\lambda\,H^{3}(\fg)[[\lambda]]$ with
$\lambda=k+\dvee$; Whitehead's lemma makes $H^{3}(\fg)$ one-dimensional
for simple $\fg$, formality of the $E_3$-operad converts both sides
into $P_3$-deformation problems, and the two Maurer--Cartan classes
agree because both are pinned by the binary bracket above. The
identification of the two formal families is the precise statement
that the derived centre of the affine algebra is the observable
algebra of quantum Chern--Simons theory on the formal disc; the
representative depends on the associator, through the same torsor as
everything else in Section~\ref{sec:ordered}.

\subsection{The operadic circle}

The tower of structures closes into a circle:
\begin{equation}
E_3\;\longrightarrow\;E_2\;\longrightarrow\;E_1
\;\longrightarrow\;E_2\;\longrightarrow\;E_3,
\end{equation}
the topological bulk determining the braided category of its lines,
the braided category its chiral boundary algebras, the boundary
algebra its ordered bar complex, the bar complex the Drinfeld double,
and the double --- through its Drinfeld centre --- the bulk again.
For the Heisenberg algebra the circle closes as a theorem: the
Drinfeld centre of the completed double-module category and the
derived chiral centre are both three-dimensional with trivial
$P_3$-structure, identified by module--comodule duality composed with
Hochschild--coHochschild duality. On the whole Koszul locus the
Morita comparison closes; the closing of the circle for the
gravitational classes is the categorical form of holography, and it
is the natural endpoint of the programme of Section~\ref{sec:bulk}.

\section{The Calabi--Yau source}\label{sec:CY}

\subsection{The Atiyah--Connes class}

A Calabi--Yau category carries a canonical obstruction to chiral
specialisation, and it is a characteristic class. Let $X$ be a smooth
projective Calabi--Yau threefold with its cyclic dg-enhancement: the
Connes operator $\Delta_2$ transports the cyclic differential to the
cochain level, the ternary product $m_3$ realises the pentagonal cell
of the Deligne structure, and their commutator is a closed
four-cochain whose Hochschild--Kostant--Rosenberg projection,
contracted with the holomorphic volume form, is
\begin{equation}
\alpha_X\;=\;\iota_{\Omega_X}\,\pi^{2,2}_{\mathrm{HKR}}
\bigl[m_3,\Delta_2\bigr]\;\in\;H^{1,2}(X),
\end{equation}
represented in Dolbeault cohomology by
$\iota_{\Omega_X}\operatorname{tr}(\bar\partial\nabla\wedge
\bar\partial\nabla)$ for any connection, independently of the
choice: the Atiyah class of $X$, curled by the cyclic structure.

\begin{theorem}[Spine]\label{thm:spine}
The two-stage specialisation --- integrate over a surface, restrict
to a curve --- carries the Atiyah--Connes class to the first
bar--cobar obstruction of the resulting chiral algebra:
$sp_{\Sigma,C}(\alpha_X)=o^{(0)}_3(A_C)$. In particular the vanishing
of $\alpha_X$ forces the vanishing of the first obstruction of every
chiral algebra the threefold produces.
\end{theorem}

\begin{proof}
The two-stage specialisation is a filtered morphism of cyclic
deformation complexes, and both sides of the equation are its
degree-two shadow: integration over $\Sigma$ carries the
Hochschild--Kostant--Rosenberg $(2,2)$-projection to the binary
collision residue, restriction to $C$ evaluates it, and the class
$[m_3,\Delta_2]$ maps to the first obstruction cocycle by naturality
of the bar differential under both operations. On $S\times E$ the
evaluation is one pairing: both sides land on the line
$H^{2}(S,\mathcal{O}_S)\otimes H^{0}(E,\Omega^{1}_E)$, and the
morphism identifies the two coefficients.
\end{proof}

The spine is the bridge between the geometric source and the chiral
theory of this paper: an obstruction living on a threefold, measured
by a Hodge class, becomes an obstruction living on a curve, measured
by the tower of Section~\ref{sec:shadow}. On a product $S\times E$
of a $K3$ surface and an elliptic curve the class concentrates, by
K\"unneth naturality, on the one-dimensional line
$H^{2}(S,\mathcal{O}_S)\otimes H^{0}(E,\Omega^{1}_E)$: the entire
question of specialisation is one scalar, and its expected value is
proportional to $\int_S c_2=24$.

\subsection{The chirality ladder}

The chiral output of a Calabi--Yau category is graded by dimension:
$d=1$ produces $E_\infty$-chiral algebras, $d=2$ produces
$E_2$-structures through the Drinfeld centre, and $d\geq3$
stabilises at $E_1$ --- the ordered theory of
Section~\ref{sec:ordered} is the generic output of the geometric
source, which is the deepest reason the ordered theory is primary.
The scalar shadow stratifies with the Hodge theory: at odd $d$ the
Serre pairing cancels the supertrace in pairs and the chiral scalar
vanishes; at even $d$ the middle Hodge number survives, giving
$\kappa_{\mathrm{ch}}(K3)=2$ and its analogues. The cohomological
Hall algebra of the resolved geometry supplies the positive half of
the quantum group, its Drinfeld double the whole; for the $K3$
surface the resulting object has twenty-four generators indexed by
the Mukai lattice, with the abelian Drinfeld coproduct
$\Delta_z x^{\pm}_i(u)=x^{\pm}_i(u)\otimes1+1\otimes
x^{\pm}_i(u-z)$, Serre relations governed by the $24\times24$
intersection matrix, and operator products computed to the order the Serre relations
require --- the quantum group of a surface, with the surface's
arithmetic in its Cartan matrix.

\subsection{One quantum group}

The families of Section~\ref{sec:ordered} and Section~\ref{sec:DS}
are specialisations of a single object $Q^{k,f,\mu}_{\fg}$,
parametrised by a Lie type, a level, a nilpotent, and a spectral
parameter: the Yangian at $k=0$, $f=0$; the affine Yangian at generic
level; the shifted Yangians by change of oscillator variables; the
Coulomb-branch algebras of three-dimensional gauge theories; and the
$\cW$-algebras, principal and not, through Drinfeld--Sokolov
reduction. One object, eight families, one duality: the inversion
principle of Theorem~\ref{thm:inversion-principle} acts on
$Q^{k,f,\mu}_{\fg}$ and restricts to every specialisation.

\section{Lattices, moonshine, and the Borcherds boundary}
\label{sec:moonshine}

\subsection{Self-duality and the census}

For a lattice algebra $V_\Lambda$ with $2$-cocycle $\varepsilon$, the
inversion principle of Theorem~\ref{thm:inversion-principle} gives
$(V_\Lambda^{\varepsilon})^{i}=(V_\Lambda^{\varepsilon^{-1}})^{c}$; for
even unimodular $\Lambda$ with symmetric cocycle the bar coalgebra is
self-dual. The lattice algebras are also the one place where a
genuinely $E_1$ chiral algebra arises classically: a non-symmetric
cocycle produces a braiding that is an independent input, not derived
from a vertex algebra. The holomorphic $c=24$ census organises itself
by the weight-one threshold: the Schellekens list \cite{Schellekens} decomposes as
$71=24+1+46$ --- twenty-four Niemeier lattice algebras, the Moonshine
module as the unique entry whose weight-one space vanishes, and forty-six
orbifolds --- and the classification data of
Section~\ref{sec:shadow} refine the census: all lattice entries are
class $\mathsf{G}$ with $\kappa=24$, the Moonshine module is class
$\mathsf{M}$ with $\kappa=12$.

\subsection{The Borcherds boundary}

The strictification mechanism of Section~\ref{sec:ordered} --- the
obstruction to a strict quantum-group structure lives in root spaces
and dies when all root multiplicities are one --- fails exactly for
Borcherds algebras, whose imaginary simple roots have multiplicity.
This is the method locating the automorphic boundary of the theory. On that boundary sit the
denominator identities: the Igusa cusp form factors as
$\Phi_{10}=\Delta_5^{2}$, with $\Delta_5$ the Borcherds product of
weight $5=c_{\Delta_5}(0)/2$ attached to the $K3$ elliptic genus by
Gritsenko--Nikulin \cite{GN}, and the second-quantised $K3$ elliptic
genus of Dijkgraaf--Moore--Verlinde--Verlinde \cite{DMVV} is the inverse of the
Igusa form,
\begin{equation}
\sum_{N\geq0}p^{N}\chi_{\mathrm{ell}}
\bigl(\operatorname{Sym}^{N}K3\bigr)
=\frac{p\,q\,y}{\Phi_{10}(p,q,y)},
\end{equation}
the symmetric-orbifold linearity of Section~\ref{sec:landscape}
exponentiated into a Borcherds product. The
weight formula $\kappa_{\mathrm{BKM}}(\Phi_N)=c_N(0)/2$ holds across
the paramodular family, with the weights $5,3,2,\tfrac32,1$ at
$N=1,2,3,4,6$: one arithmetic law, five instances, all Borcherds
products. Against this proved automorphic arithmetic
stands Conjecture~\ref{conj:mukai}: that the chiral algebra recognised
from the $K3$ source carries conductor $2c_+=8$ and quantisation
$\hbar^{2}=-1/8$. The five scalars attached to $K3\times E$ ---
categorical $0$, Hodge $0$, Heisenberg $3$, Borcherds weight $5$,
fibre Euler number $24$ --- are five distinct invariants, and the
conjecture identifies exactly one of them with the conductor. An order-two datum makes the
distinction vivid: the physical partition function is
$\Delta_5^{-2}$, and the orientation bit that distinguishes
the two candidate chiral objects is remembered by the algebra and
integrated away by its partition function --- the sharpest instance of
the principle that the object carries strictly more than its scalar
shadow.

\subsection{What the tower sees of moonshine}

The shadow tower sees the Monster through its coefficients: for
$V^{\natural}$, $\Delta=20/71$ and the tower is asymptotic with the
same Borel structure as Ising; the McKay--Thompson refinements enter
as twining fingerprints; and moonshine enters the classification as
boundary data, the automorphic side of the correspondence standing on
its own theorems.

\section{The physics}\label{sec:physics}

Each statement of this section is a theorem of the formalism; where a
comparison map with a physical theory enters, it enters in the
statement.

\subsection{The anomaly is the curvature}

\begin{theorem}[Rigidity of the anomaly]\label{thm:anomaly}
The obstruction to a translation-invariant chiral coproduct on
$\Vir_c$ is the class $(c/2)\,\Omega^{\mathrm{univ}}
\in H^{2}_{\mathrm{ch}}(\Vir_c,\Vir_c)$: a $z$-independent coproduct
exists if and only if $c=0$, and the coefficient $c/2=\kappa(\Vir_c)$
is stable under pullback along modules. The conformal anomaly of
two-dimensional physics, the fourth-order pole of the stress tensor,
and the modular curvature of Theorem~\ref{thm:curvature} are one
number.
\end{theorem}

The genus-one free energy is triply determined: as the degree-two
projection of $\Theta_{\cA}$, as the logarithmic derivative
$\partial_\tau\log\eta^{-2\kappa}$, and as the regularised Casimir
energy on the circle --- three computations, one number,
$F_1=\kappa/24$. The string-theoretic identities follow at scalar
level. The ghost
system carries $\kappa(bc_{\lambda=2})=-13$; matter plus ghosts is
curvature-free precisely when $c=26$, which is the identity
$\varkappa(\Vir)=13$ of Theorem~\ref{thm:complementarity} read as
anomaly cancellation, with the self-dual point $c=13$ at its centre.
The genus-zero bar complex of the $c=26$ system is identified with the
BRST complex under a filtered comparison; nilpotence of the BRST
charge, $Q^{2}=\frac{c-26}{12}\sum_n(n^{3}-n):\!c_{-n}c_n\!:$, is the
uncurved condition.

\subsection{The two central charges}

The theory distinguishes two exceptional values of $c$ and assigns
each its own mechanism. At $c=26$ the \emph{dual} curvature vanishes:
$\kappa(\Vir_{26-c})=0$, matter and ghosts cancel, and the total bar
differential of the coupled system is uncurved --- the critical
dimension is the flatness locus of the Koszul-dual theory. At $c=13$
the theory is \emph{self-dual}: the fixed point of the involution,
the balance point of the entropy sum rule, the centre of the
conductor. Both values are visible already in
\eqref{eq:conductor-values}, as the zero and the fixed point of one
involution; the string selects the first, the symmetry the second.

\subsection{The anomaly ratio}

The ratio $\varrho=\kappa/c=\sum_i 1/(m_i+1)$, summed over the
exponents of the underlying Lie algebra, is a spectral invariant of
the family: $1$ for the Heisenberg lane, $\tfrac12$ for Virasoro,
$\tfrac56$ for $\cW_3$, $\tfrac{13}{12}$ for $\cW_4$,
$\tfrac{121}{126}$ for the $E_8$ family. It converts between the two
conductors, $\varkappa=\varrho\,K$, and it measures what survives
Drinfeld--Sokolov reduction: the exponents are the surviving weights,
and the ratio is their harmonic signature. That a single rational
number carries the exponent spectrum of the reduction into the
modular curvature is the kind of bookkeeping identity the subject
produces when the definitions are right.

\subsection{The string measure}

For the rank-$r$ Heisenberg algebra the analytic theory closes. The
sewing operator on the weighted Fock space is trace class with
eigenvalues $q^{n}$; its Fredholm determinant is
$\det(1-T_q)=q^{-1/24}\eta(q)$; and the genus-$g$ partition function,
assembled by sewing, is
\begin{equation}
Z_g\;=\;\bigl(\det\operatorname{Im}\Omega\bigr)^{-r/2}
\bigl({\det}'_{\zeta}\Delta_{\Sigma_g}\bigr)^{-r/2},
\end{equation}
the Polyakov measure of the bosonic string \cite{DP}, derived here by
operator theory from the bar complex, with the Belavin--Knizhnik form
as cross-check. Lattice algebras multiply this by the Siegel theta
series; at genus two the theta series separates $E_8\oplus E_8$ from
$D_{16}^{+}$, a separation genus two alone achieves.

\subsection{Soft limits and primitivity}

The first shadows have names in the flat-space expansion of gravity:
$S_2=c/2$ is the leading (Weinberg) soft coefficient, $S_3=2$ the
subleading (Cachazo--Strominger) coefficient, and
$S_4=10/[c(5c+22)]$ the first coefficient beyond, entering through the
Mellin residue at $\Delta_s=4-r$; the higher soft expansion proceeds
through the Mellin residues, each shadow coefficient feeding the pole
at its own depth. A cleaner
statement is coalgebraic: for every class-$\mathsf{M}$ algebra obtained
by Drinfeld--Sokolov reduction, the coproduct on the dual is strictly
primitive --- a ghost-number obstruction kills all group-like
corrections --- while the $A_\infty$ tower is infinite. Gauge theory
deforms the coproduct and truncates the tower; gravity does the
opposite. Gravity braids particles; braiding is its entire product.

\subsection{Three-dimensional gravity}

The Virasoro instance of Theorem~\ref{thm:terminality} packages the
boundary algebra at Brown--Henneaux central charge $c=3\ell/2G$, the
bulk $\cZ^{\mathrm{der}}_{\mathrm{ch}}(\Vir_c)$ with Poincar\'e
polynomial $1+t^{2}+t^{3}$, and the collision kernel
$r_c(z)=(c/2)z^{-3}+2Tz^{-1}$ into an algebraic holographic sector for
three-dimensional gravity. Within it: BTZ states are Maurer--Cartan
elements; the Cardy entropy is the character asymptotics under modular
invariance and vacuum dominance; and the two-branch entropy shadow
crosses at $t_{*}=3S_{\mathrm{BH}}/13$, the coefficient being
$c$-independent because the rate sum is $\varkappa/3=13/3$. In the Cardy regime the entropy is character asymptotics: given
modular invariance of the torus block and dominance of the vacuum
channel, the density of states at large weight gives
$S=2\pi\sqrt{c\Delta/6}+2\pi\sqrt{c\bar\Delta/6}$, with the
determinant of the modular transformation supplying the
$-\tfrac32\log S$ correction; for the Brown--Henneaux value of $c$
this is the horizon area over $4G$. The
crossing formula is a theorem; the two-branch structure realises
Conjecture~\ref{conj:reflection}, whose scalar shadow
$\varkappa=13$ it uses; and the dynamical-metric completion of the
sector is the frontier of the subject, with the sum over topologies
replaced by the Maurer--Cartan induction of Theorem~\ref{thm:mc} ---
an induction, every term finite.

\subsection{Chern--Simons theory and knots}

The abelian case anchors the dictionary at the level of theorems: the
Heisenberg algebra $\mathcal{H}_k$ is the boundary algebra of abelian
Chern--Simons theory at level $k$, the boundary operator product
$k/(z-w)^{2}$ arising from the slab Green's function, and
$\kappa=k$ identifies the modular curvature with the Chern--Simons
level; the Wilson-line framing phase $e^{2\pi iq^{2}/k}$ is the
boundary value of the $R$-matrix $e^{k\hbar/z}$. In the nonabelian
case the ordered bar complex of $\widehat{\fg}_k$ carries the braid
monodromy of the quantum group at $q=e^{\pi i/(k+\dvee)}$; closure of
braids through the genus-one trace, together with the standard
ribbon-category trace, computes the coloured Jones polynomials, the
writhe correction being the Casimir framing anomaly. Level--rank and
boundary dualities are instances of Koszul duality of boundary
conditions; the identification of the Yangian with the line-operator
algebra of four-dimensional Chern--Simons theory \cite{CWY} is the
physical face of Theorem~\ref{thm:yangian}, with the
non-renormalisation of the line OPE as the physical face of
Koszulness.

\subsection{Wilson lines and evaluation modules}

Line operators of the three-dimensional theory are evaluation modules
of the ordered dual: the spectral parameter is the holomorphic
position of the line, braiding of lines is the $R$-matrix, and the
operator product of parallel lines is one-loop exact --- the
non-renormalisation theorem of the physics is the collapse of the bar
spectral sequence, which is Koszulness. Monopole operators arrive
with the affine Grassmannian, and the Coulomb branches of
three-dimensional gauge theories realise the shifted Yangians \cite{BFN}: the
ordered theory of Section~\ref{sec:ordered} is the algebra of the
line-operator sector, stated once and specialised everywhere.

\subsection{Entanglement}

The genus-one shadow assigns to an interval the entropy
$S_{EE}=(c/3)\log(L/\varepsilon)$, and complementarity of
Theorem~\ref{thm:complementarity} gives the sum rule
$S_{EE}(\Vir_c)+S_{EE}(\Vir_{26-c})=\tfrac{26}{3}\log(L/\varepsilon)$,
saturated symmetrically at $c=13$. The classification grades
entanglement complexity: Gaussian for class $\mathsf{G}$, a single
logarithm for $\mathsf{L}$, a quartic contact correction for
$\mathsf{C}$, an infinite R\'enyi tower for $\mathsf{M}$. The
Lagrangian splitting of Theorem~\ref{thm:complementarity} is a
symplectic code --- each summand isotropic, the cross-pairing perfect
--- and a Hilbert-space datum completes it to a quantum code; the
formulas of this subsection are shadow functionals, exact as stated.

\subsection{Integrable structures}

The genus-zero shadow connection of the affine algebra is the
Knizhnik--Zamolodchikov connection with coupling $1/(k+\dvee)$;
flatness is the Arnold relation contracted with the Casimir, which is
the classical Yang--Baxter equation; the Gaudin Hamiltonians are its
residues. At genus one the propagator degenerates through the
Weierstrass zeta function to the elliptic $r$-matrices of Belavin and
the dynamical connection of Knizhnik--Zamolodchikov--Bernard and Felder \cite{Felder}, and the
$B$-cycle monodromy exponentiates the coupling to
$q=e^{\pi i/(k+\dvee)}$; braid monodromy of the ordered bar complex
computes the coloured Jones polynomials after the standard
ribbon-category trace. All three appearances of $\hbar$ --- the
KZ coupling, the line-operator loop parameter, the bar coupling ---
are the same number $1/(k+\dvee)$, and this identification is forced
by the Sugawara construction.

\section{The one question}\label{sec:question}

Every construction in this paper is a projection of the single element
$\Theta_{\cA}\in\MC(\gmod_{\cA})$: the $r$-matrix and the
Knizhnik--Zamolodchikov connection at genus zero; $\kappa$ and the
curvature at degree two; the shadow tower along a line; the
complementarity splitting and its conductor over $\Mbar_g$; the genus
expansion as a trace; the Verlinde numbers on the integrable quotient;
the quantum group in the ordered kernel. Everything proved is a
finite-order projection of $\Theta_{\cA}$. Everything open is a
question about its ambient.

\begin{question}
Characterise the modular convolution algebra $\gmod_{\cA}$ as a
geometric object.
\end{question}

The evidence assembled here proposes an answer: $\gmod_{\cA}$ is the
endomorphism algebra of a Lagrangian embedding
$\cL_{\cA}\hookrightarrow\cM_{\mathrm{vac}}(\cA)$ in a
$(-2)$-shifted symplectic stack of vacua. The shadow tower computes
its Taylor jets; the Swiss-cheese structure of
Section~\ref{sec:bulk} its local composition law; the shift matches
the $-(3g-3)$-shifted symplectic structure of
Theorem~\ref{thm:complementarity}. The remaining object is the global
geometry: the stack itself, the moduli of its Lagrangians, and the
descent that makes the construction independent of all choices.
The known boundary of the theory --- the critical level, where
$\kappa=0$ and the centre becomes the algebra of functions on opers;
the logarithmic algebras outside the classification; the
multi-weight corrections that leave $\bQ(c)$; the Borcherds algebras
where the strictification mechanism fails and the automorphic forms
begin --- marks exactly where that geometry must be nontrivial.

Five problems structure the frontier, and each is a construction.

\emph{Chain-level topologisation.} Lift the $E_3$-structure of
Theorem~\ref{thm:E3} from cohomology to the original complex for the
gravitational classes: the coherent tower over $[m,G]=\partial_z$,
expected to close in the coderived ambient, where the obstruction
cocycles are curvature-torsion.

\emph{Nodal descent.} Extend Theorems~\ref{thm:pbw}
and~\ref{thm:complementarity} across the boundary of $\Mbar_{g,n}$
with prescribed clutching: logarithmic sewing data at the separating
and non-separating divisors, completing the modular functor at chain
level.

\emph{The multi-weight generating function.} A closed form for the
cross-channel corrections of Theorem~\ref{thm:multiweight} in the
Frobenius geometry of the $\cW_N$ families, where the computed
values at genus two, three and four already fix the shape.

\emph{Compact completion.} Extend the ordered duality of
Theorem~\ref{thm:yangian} from the evaluation core to the full
module category: in type $A$ a single compact-generation statement
on the Baxter locus.

\emph{The dynamical metric.} Complete the algebraic holographic
sector of Section~\ref{sec:physics} to a gravitational path
integral, with the Maurer--Cartan induction as the sum over
topologies and the Igusa kernel as the genus-two input.

One $1$-form, one relation, one extraction, one ordering.

\appendix

\section{The landscape census}\label{app:census}

\begin{center}\small
\begin{tabular}{lllll}
\emph{algebra} & \emph{class} & $\kappa$ & \emph{dual partner} &
$\varkappa$\\
\hline
$\mathcal{H}_k$ & $\mathsf{G}$ & $k$ &
curved $\operatorname{Sym}$ & $0$\\
free fermion & $\mathsf{G}$ & $1/4$ & itself & --- \\
$V_\Lambda$, $\Lambda$ unimodular & $\mathsf{G}$ &
$\operatorname{rk}\Lambda$ & itself (cocycle inverted) & $0$\\
$V_k(\fg)$, $k\neq-\dvee$ & $\mathsf{L}$ &
$\frac{(k+\dvee)\dim\fg}{2\dvee}$ &
level $-k-2\dvee$ & $0$\\
$\widehat{\fg}_1$, simply laced & $\mathsf{G}$ &
$\operatorname{rk}\fg$ & lattice realisation & $0$\\
$\beta\gamma_\lambda$ & $\mathsf{C}$ & $6\lambda^{2}-6\lambda+1$ &
$bc_\lambda$ & $0$\\
$\Vir_c$ & $\mathsf{M}$ & $c/2$ &
$\Vir_{26-c}$ (Conj.~\ref{conj:reflection}) & $13$\\
$\cW_N$ & $\mathsf{M}$ & $c(H_N-1)$ &
level-reflected $\cW_N$ & $(H_N{-}1)K_N$\\
$\cW^{k}(\mathfrak{sl}_3,f_{\min})$ & mixed & see \S\ref{sec:DS} &
$k\mapsto-k-6$ & $c$-conductor $50$\\
$V^{\natural}$ & $\mathsf{M}$ & $12$ & Virasoro sector partner & ---\\
Leech $V_\Lambda$ & $\mathsf{G}$ & $24$ & itself & $0$\\
$Y_\hbar(\mathfrak{sl}_2)$ & $E_1$ & --- &
$Y_{-\hbar}(\mathfrak{sl}_2)$ & ---\\
$V_{-\dvee}(\fg)$ & $\mathsf{FF}$ & $0$ &
$\mathrm{Fun}\,\mathrm{Op}_{{}^{L}\fg}(D)$ & $0$
\end{tabular}
\end{center}

\medskip
\noindent
The classes are those of Theorem~\ref{thm:classification}; the
$\cW_N$ conductor carries the Drinfeld--Sokolov transport hypothesis of
Theorem~\ref{thm:complementarity-values}; the central-charge
conductors $c+c^{!}$ are $26$, $100$, $4N^{3}-2N-2$, $2\dim\fg$, $50$
on their respective rows, with $K_N=4N^{3}-2N-2$, exact as rational
functions of the parameter.

\section{Notation}\label{app:notation}

\begin{center}
\begin{tabular}{ll}
$\eta_{ij}=d\log(z_i-z_j)$ & the logarithmic propagator\\
$\FM_n(X)$ & Fulton--MacPherson compactification of $n$ points\\
$\barB^{\mathrm{ord}}_X(\cA)$, $\barB^{\Sigma}_X(\cA)$ &
ordered and symmetric bar coalgebras\\
$\cA^{i}$, $\cA^{!}$ & quadratic dual coalgebra; Koszul dual algebra\\
$K_X(\cA)=\mathbb{D}_{\Ran}\barB_X(\cA)$ & Verdier algebra of the bar
coalgebra\\
$\cZ^{\mathrm{der}}_{\mathrm{ch}}(\cA)$ & derived chiral centre
(chiral Hochschild complex)\\
$\Theta_{\cA}=D_{\cA}-d_0$ & the modular Maurer--Cartan element\\
$\gmod_{\cA}$ & modular convolution Lie algebra\\
$\kappa(\cA)$ & modular characteristic: scalar of the genus curvature\\
$S_r$, $H(t)$, $Q(t)$ & shadow tower, its generating function, its
square\\
$\Delta=8\kappa S_4$ & critical discriminant of the tower\\
$r_{\max}$ & shadow depth; classes $\mathsf{G},\mathsf{L},
\mathsf{C},\mathsf{M},\mathsf{FF}$\\
$\varkappa=\kappa+\kappa^{!}$ & conductor; $K=c+c^{!}$ its
central-charge form\\
$\lambda_g^{\mathrm{FP}}$ & Faber--Pandharipande Hodge numbers\\
$r(z)$, $R(z)$ & classical and quantum $r$-matrices\\
$E(z,w)$ & the prime form; $\omega_g=c_1(\bE_g)$ the Hodge class\\
$q=e^{\pi i/(k+\dvee)}$ & the quantum parameter of level $k$
\end{tabular}
\end{center}

\begin{thebibliography}{99}
\bibitem[Ara]{Arakawa} T.~Arakawa,
\emph{Representation theory of $W$-algebras},
Invent. Math. \textbf{169} (2007), 219--320.
\bibitem[Ar]{Arnold} V.~I.~Arnold,
\emph{The cohomology ring of the colored braid group},
Mat. Zametki \textbf{5} (1969), 227--231.
\bibitem[BDSK]{BDSK} B.~Bakalov, A.~De Sole, V.~Kac,
\emph{Computation of cohomology of vertex algebras},
Jpn. J. Math. \textbf{16} (2021), 81--154.
\bibitem[BV]{BV} D.~Barbasch, D.~Vogan,
\emph{Unipotent representations of complex semisimple groups},
Ann. of Math. \textbf{121} (1985), 41--110.
\bibitem[BD]{BD} A.~Beilinson, V.~Drinfeld,
\emph{Chiral Algebras},
AMS Colloquium Publications \textbf{51}, 2004.
\bibitem[BK]{BK} A.~Belavin, V.~Knizhnik,
\emph{Algebraic geometry and the geometry of quantum strings},
Phys. Lett. B \textbf{168} (1986), 201--206.
\bibitem[Zhu]{Zhu} Y.~Zhu,
\emph{Modular invariance of characters of vertex operator algebras},
J. Amer. Math. Soc. \textbf{9} (1996), 237--302.
\bibitem[Bo]{Borcherds} R.~Borcherds,
\emph{Vertex algebras, Kac--Moody algebras, and the Monster},
Proc. Natl. Acad. Sci. USA \textbf{83} (1986), 3068--3071.
\bibitem[BH]{BH} J.~D.~Brown, M.~Henneaux,
\emph{Central charges in the canonical realization of asymptotic
symmetries},
Comm. Math. Phys. \textbf{104} (1986), 207--226.
\bibitem[CG]{CG} K.~Costello, O.~Gwilliam,
\emph{Factorization Algebras in Quantum Field Theory}, I--II,
Cambridge University Press, 2017, 2021.
\bibitem[CWY]{CWY} K.~Costello, E.~Witten, M.~Yamazaki,
\emph{Gauge theory and integrability, I--II},
ICCM Not. \textbf{6} (2018).
\bibitem[DMVV]{DMVV} R.~Dijkgraaf, G.~Moore, E.~Verlinde,
H.~Verlinde,
\emph{Elliptic genera of symmetric products and second quantized
strings},
Comm. Math. Phys. \textbf{185} (1997), 197--209.
\bibitem[DP]{DP} E.~D'Hoker, D.~H.~Phong,
\emph{The geometry of string perturbation theory},
Rev. Mod. Phys. \textbf{60} (1988), 917--1065.
\bibitem[Dr]{Drinfeld} V.~Drinfeld,
\emph{Quantum groups},
Proc. ICM Berkeley (1986), 798--820.
\bibitem[EK]{EK} P.~Etingof, D.~Kazhdan,
\emph{Quantization of Lie bialgebras}, I--V,
Selecta Math. (1996--2000).
\bibitem[FP]{FP} C.~Faber, R.~Pandharipande,
\emph{Hodge integrals and Gromov--Witten theory},
Invent. Math. \textbf{139} (2000), 173--199.
\bibitem[Fay]{Fay} J.~Fay,
\emph{Theta Functions on Riemann Surfaces},
Lecture Notes in Math. \textbf{352}, Springer, 1973.
\bibitem[Fel]{Felder} G.~Felder,
\emph{Conformal field theory and integrable systems associated to
elliptic curves},
Proc. ICM Z\"urich (1994), 1247--1255.
\bibitem[FF]{FF} B.~Feigin, E.~Frenkel,
\emph{Affine Kac--Moody algebras at the critical level and
Gelfand--Dikii algebras},
Int. J. Mod. Phys. A \textbf{7} (1992), 197--215.
\bibitem[FG]{FG} J.~Francis, D.~Gaitsgory,
\emph{Chiral Koszul duality},
Selecta Math. \textbf{18} (2012), 27--87.
\bibitem[FH]{FH} E.~Frenkel, D.~Hernandez,
\emph{Baxter's relations and spectra of quantum integrable models},
Duke Math. J. \textbf{164} (2015), 2407--2460.
\bibitem[FM]{FM} W.~Fulton, R.~MacPherson,
\emph{A compactification of configuration spaces},
Ann. of Math. \textbf{139} (1994), 183--225.
\bibitem[GK]{GK} E.~Getzler, M.~Kapranov,
\emph{Modular operads},
Compositio Math. \textbf{110} (1998), 65--126.
\bibitem[GTL]{GTL} S.~Gautam, V.~Toledano Laredo,
\emph{Yangians and quantum loop algebras},
Selecta Math. \textbf{19} (2013), 271--336.
\bibitem[HJ]{HJ} D.~Hernandez, M.~Jimbo,
\emph{Asymptotic representations and Drinfeld rational fractions},
Compositio Math. \textbf{148} (2012), 1593--1623.
\bibitem[HZ]{HZ} J.~Harer, D.~Zagier,
\emph{The Euler characteristic of the moduli space of curves},
Invent. Math. \textbf{85} (1986), 457--485.
\bibitem[GN]{GN} V.~Gritsenko, V.~Nikulin,
\emph{Siegel automorphic form corrections of some Lorentzian
Kac--Moody Lie algebras},
Amer. J. Math. \textbf{119} (1997), 181--224.
\bibitem[KW]{KW} V.~Kac, M.~Wakimoto,
\emph{Modular invariant representations of infinite-dimensional Lie
algebras and superalgebras},
Proc. Natl. Acad. Sci. USA \textbf{85} (1988), 4956--4960.
\bibitem[KZ]{KZ} V.~Knizhnik, A.~Zamolodchikov,
\emph{Current algebra and Wess--Zumino model in two dimensions},
Nucl. Phys. B \textbf{247} (1984), 83--103.
\bibitem[Ko]{Kontsevich} M.~Kontsevich,
\emph{Deformation quantization of Poisson manifolds},
Lett. Math. Phys. \textbf{66} (2003), 157--216.
\bibitem[Lor]{Vol1} R.~Lorgat,
\emph{Modular Homotopy Theory for Algebraic Factorization Algebras
on Algebraic Curves, Vol.~I: $E_1$ Chiral Algebras and Modular Koszul
Duality}, monograph, 2026.
\bibitem[LV]{LV} J.-L.~Loday, B.~Vallette,
\emph{Algebraic Operads},
Grundlehren \textbf{346}, Springer, 2012.
\bibitem[MO]{MO} D.~Maulik, A.~Okounkov,
\emph{Quantum groups and quantum cohomology},
Ast\'erisque \textbf{408} (2019).
\bibitem[Mum]{Mum} D.~Mumford,
\emph{Towards an enumerative geometry of the moduli space of curves},
in \emph{Arithmetic and Geometry II}, Birkh\"auser, 1983, 271--328.
\bibitem[Pol]{Polyakov} A.~M.~Polyakov,
\emph{Quantum geometry of bosonic strings},
Phys. Lett. B \textbf{103} (1981), 207--210.
\bibitem[Pos]{Pos} L.~Positselski,
\emph{Two kinds of derived categories, Koszul duality, and
comodule-contramodule correspondence},
Mem. Amer. Math. Soc. \textbf{212} (2011).
\bibitem[BFN]{BFN} A.~Braverman, M.~Finkelberg, H.~Nakajima,
\emph{Coulomb branches of $3d$ $\mathcal{N}=4$ quiver gauge theories
and slices in the affine Grassmannian},
Adv. Theor. Math. Phys. \textbf{23} (2019), 75--166.
\bibitem[Pr]{Priddy} S.~Priddy,
\emph{Koszul resolutions},
Trans. Amer. Math. Soc. \textbf{152} (1970), 39--60.
\bibitem[Sch]{Schellekens} A.~N.~Schellekens,
\emph{Meromorphic $c=24$ conformal field theories},
Comm. Math. Phys. \textbf{153} (1993), 159--185.
\bibitem[Ta]{Tamarkin} D.~Tamarkin,
\emph{Another proof of M. Kontsevich formality theorem},
arXiv:math/9803025.
\bibitem[TUY]{TUY} A.~Tsuchiya, K.~Ueno, Y.~Yamada,
\emph{Conformal field theory on universal family of stable curves
with gauge symmetries},
Adv. Stud. Pure Math. \textbf{19} (1989), 459--566.
\bibitem[Ver]{Verlinde} E.~Verlinde,
\emph{Fusion rules and modular transformations in 2D conformal field
theory},
Nucl. Phys. B \textbf{300} (1988), 360--376.
\bibitem[Wak]{Wakimoto} M.~Wakimoto,
\emph{Fock representations of the affine Lie algebra
$A^{(1)}_1$},
Comm. Math. Phys. \textbf{104} (1986), 605--609.
\bibitem[Wit]{WittenVol} E.~Witten,
\emph{On quantum gauge theories in two dimensions},
Comm. Math. Phys. \textbf{141} (1991), 153--209.
\bibitem[Zam]{Zam} A.~B.~Zamolodchikov,
\emph{Infinite additional symmetries in two-dimensional conformal
quantum field theory},
Theor. Math. Phys. \textbf{65} (1985), 1205--1213.
\end{thebibliography}

\end{document}
